% pyCMF 1.0 user and developer manual.
% Build:  latexmk -pdf pycmf_manual.tex   (or pdflatex twice)
\documentclass[11pt,a4paper,oneside]{report}

\usepackage[T1]{fontenc}
\usepackage[utf8]{inputenc}
\usepackage{lmodern}
\usepackage{textcomp}
\usepackage{microtype}
\usepackage[margin=2.5cm,headheight=14pt]{geometry}
\usepackage{amsmath,amssymb}
\usepackage{booktabs}
\usepackage{tabularx}
\usepackage{longtable}
\usepackage{array}
\usepackage{enumitem}
\usepackage{listings}
\usepackage{fancyhdr}
\usepackage{titlesec}
\usepackage[hidelinks]{hyperref}

% ---------------------------------------------------------------- layout
\setlength{\parindent}{0pt}
\setlength{\parskip}{0.6em plus 0.2em minus 0.1em}
\setlist{itemsep=0.2em, topsep=0.3em}
\renewcommand{\arraystretch}{1.15}
\setlength{\emergencystretch}{3em}

\titleformat{\chapter}[hang]{\normalfont\huge\bfseries}{\thechapter}{1em}{}
\titlespacing*{\chapter}{0pt}{0pt}{2em}

\pagestyle{fancy}
\fancyhf{}
\fancyhead[L]{\small pyCMF 1.0 Manual}
\fancyhead[R]{\small\nouppercase{\leftmark}}
\fancyfoot[C]{\small\thepage}
\renewcommand{\headrulewidth}{0.4pt}
\fancypagestyle{plain}{\fancyhf{}\fancyfoot[C]{\small\thepage}\renewcommand{\headrulewidth}{0pt}}

% ---------------------------------------------------------------- code
\lstset{
  basicstyle=\ttfamily\small,
  columns=fullflexible,
  keepspaces=true,
  frame=single,
  framerule=0.4pt,
  xleftmargin=4pt,
  xrightmargin=4pt,
  framexleftmargin=2pt,
  aboveskip=0.8em,
  belowskip=0.8em,
  showstringspaces=false,
  breaklines=true,
  breakatwhitespace=true,
  upquote=true,
  tabsize=4,
}
\lstdefinestyle{python}{language=Python, morekeywords={self}}
\lstdefinestyle{shell}{language=bash}
\lstdefinestyle{output}{frame=leftline}

% inline code: characters printed literally, line breaks allowed after
% _ , . ( = so that long names fit into narrow table columns
\ExplSyntaxOn
\str_new:N \l__code_str
\NewDocumentCommand{\code}{m}
  {
    \str_set:Nn \l__code_str {#1}
    \texttt
      {
        \str_map_inline:Nn \l__code_str
          {
            \str_case:nnF {##1}
              {
                {_} {\_\allowbreak}
                {,} {,\allowbreak}
                {.} {.\allowbreak}
                {(} {(\allowbreak}
                {=} {=\allowbreak}
                {'} {\textquotesingle}
              }
              {##1}
          }
      }
  }
\ExplSyntaxOff
\newcommand{\pycmf}{pyCMF}
\newcommand{\pyscf}{PySCF}
\newcommand{\fbar}{\bar f}
\newcommand{\Nk}{N_k}

\hypersetup{
  pdftitle={pyCMF 1.0 Manual},
  pdfauthor={QuantumLab HCMUS},
  pdfsubject={Correlated mean-field methods built on PySCF},
}

\begin{document}

% ---------------------------------------------------------------- title
\begin{titlepage}
  \centering
  \vspace*{3cm}
  {\Huge\bfseries pyCMF\par}
  \vspace{0.6cm}
  {\Large Correlated Mean-Field Methods Built on PySCF\par}
  \vspace{1.5cm}
  {\Large User and Developer Manual\par}
  \vspace{0.4cm}
  {\large Version 1.0.0\par}
  \vfill
  {\large QuantumLab HCMUS\par}
  \vspace{0.2cm}
  {University of Science, Vietnam National University Ho Chi Minh City\par}
  \vspace{1cm}
  {November 2026\par}
  \vspace{0.5cm}
  {\small Source code: \url{https://github.com/QuantumLab-HCMUS/pyCMF-v1}\par}
\end{titlepage}

\pagenumbering{roman}
\tableofcontents
\clearpage
\pagenumbering{arabic}

\chapter{Introduction}
\label{ch:intro}

\section{What pyCMF is}

\pycmf{} is a Python package of \emph{correlated mean-field methods} built on
\pyscf{}. Its core method is one-body second-order M\o ller--Plesset
perturbation theory (OBMP2). In OBMP2 the Hartree--Fock Fock operator is
replaced by a \emph{correlated Fock operator} that contains a one-body MP2
correlation potential. The orbitals, the orbital energies and the amplitudes
are then iterated to self-consistency, at the cost of one MP2 calculation per
cycle.

Because the correlated Fock operator is a one-body operator, the methods of
\pycmf{} behave like a mean-field theory in many respects. They give orbitals,
orbital energies and a single determinant, they can be started from a
non-Aufbau determinant to describe excited states, and they can be used to
dress the one-body part of an active-space Hamiltonian. On top of OBMP2 the
package provides:

\begin{itemize}
  \item O2BMP2, which keeps only the opposite-spin amplitudes and scales them;
  \item one-body double hybrids (OBDH, O2BDH), in which a Kohn--Sham Fock
        matrix is augmented by a scaled OBMP2 correlation potential;
  \item state-specific excited states with the maximum overlap method (MOM)
        and its initial variant (IMOM);
  \item one-body downfolding (OBDF), which produces an active-space
        Hamiltonian for any many-body solver;
  \item properties: densities, dipole moments, $\langle S^2\rangle$ and
        isotropic hyperfine coupling constants;
  \item experimental extensions: restricted open-shell OBMP2, projection
        embedding in DFT, and OBMP2 and downfolding for periodic systems with
        $k$-point sampling.
\end{itemize}

\section{Methods and their status}

Table~\ref{tab:methods} lists the methods of version~1.0. \emph{Stable}
methods are covered by the test suite and exported at the top level of the
package (\code{import pycmf}). \emph{Experimental} methods are under
development: they are imported from their sub-package, emit an
\code{ExperimentalWarning} when they are instantiated, and are described in
Chapter~\ref{ch:validation} together with their known issues.

\begin{table}[!htb]
\centering
\small
\caption{Methods of \pycmf{} 1.0.}
\label{tab:methods}
\begin{tabularx}{\textwidth}{@{}l X l l@{}}
\toprule
Method & Description & Entry point & Status \\
\midrule
OBMP2 & self-consistent one-body MP2, restricted and unrestricted & \code{pycmf.OBMP2} & stable \\
O2BMP2 & OBMP2 with scaled opposite-spin amplitudes only & \code{pycmf.O2BMP2} & stable \\
OBDH & one-body double hybrid: Kohn--Sham Fock matrix plus scaled OBMP2 potential & \code{pycmf.OBDH} & stable \\
O2BDH & OBDH with scaled opposite-spin amplitudes only & \code{pycmf.O2BDH} & stable \\
OBDF & one-body downfolding to an active-space Hamiltonian & \code{pycmf.OBDF} & stable \\
excited states & any of the first four with \code{occupation='mom'} or \code{'imom'} & keyword & stable \\
RO-OBMP2 & restricted open-shell OBMP2 & \code{pycmf.mol.roobmp2.ROOBMP2} & experimental \\
embedding & OBMP2- and OBDH-in-DFT projection embedding, concentric localization & \code{pycmf.embed} & experimental \\
KOBMP2, KO2BMP2 & periodic systems, closed shell: density fitting, $k$-point symmetry, MPI & \code{pycmf.pbc} & experimental \\
KOBDF & downfolding to active bands with $k$-points & \code{pycmf.pbc.KOBDF} & experimental \\
KUOBMP2 & periodic systems, unrestricted (code of version 0.2) & \code{pycmf.pbc.KUOBMP2} & experimental \\
\bottomrule
\end{tabularx}
\end{table}

\section{Conventions of this manual}

\begin{itemize}
  \item Energies are in Hartree unless stated otherwise. For periodic
        systems all energies are per unit cell.
  \item Indices $i,j,k,l$ denote occupied orbitals, $a,b,c,d$ virtual
        orbitals and $p,q,r,s$ general orbitals. Two-electron integrals are
        written in chemists' notation, $(pq|rs)$.
  \item Python code is shown in framed boxes. Lines that start with
        \code{$} are shell commands.
  \item Names in \texttt{typewriter} font are Python objects or files of the
        repository. Paths are relative to the root of the repository.
\end{itemize}

\section{Design principles}

\pycmf{} follows the conventions of \pyscf{}, so that a user of \pyscf{} can
use it without learning a new interface:

\begin{enumerate}
  \item A method object is built from a converged mean-field object of
        \pyscf{}. Options are attributes of the object.
  \item \code{kernel()} (or \code{run()}, which returns the object itself)
        performs the calculation and returns the main result, the total
        energy. All other results are attributes.
  \item The input mean-field object is never modified.
  \item Each method has exactly one implementation. Density fitting, DIIS,
        the occupation scheme, the spin scaling and the order of the
        Baker--Campbell--Hausdorff expansion are attributes of that
        implementation, not separate classes or modules.
\end{enumerate}

\section{How to cite}

If you use \pycmf{}, please cite the papers of the methods you use
(Appendix~\ref{ch:references}) and \pyscf{}. The file \code{CITATION.cff} in
the repository contains the same information in machine-readable form.

\begin{itemize}
  \item OBMP2: Lan and Yanai, J.~Chem.~Phys.~\textbf{138}, 224108 (2013);
        Tran, J.~Phys.~Chem.~A \textbf{125}, 9242 (2021).
  \item Densities: Tran, Phys.~Chem.~Chem.~Phys.~\textbf{24}, 19393 (2022).
  \item O2BMP2: Tran, Nguyen and Tran, J.~Phys.~Chem.~A \textbf{128}, 1543 (2024).
  \item State-specific excited states: Tran and Tran, J.~Chem.~Phys.~\textbf{162}, 104109 (2025).
  \item Downfolding: Le and Tran, J.~Phys.~Chem.~A \textbf{127}, 5222 (2023).
  \item \pyscf{}: Sun \emph{et al.}, WIREs Comput.~Mol.~Sci.~\textbf{8}, e1340 (2018).
\end{itemize}

\section{License}

\pycmf{} is distributed under the Apache License, Version~2.0 (files
\code{LICENSE} and \code{NOTICE}). It is developed by QuantumLab HCMUS,
University of Science, Vietnam National University Ho Chi Minh City.

\chapter{Installation}
\label{ch:install}

\section{Requirements}

\begin{itemize}
  \item Python 3.9 or newer;
  \item NumPy 1.22 or newer and SciPy 1.8 or newer;
  \item \pyscf{} 2.10 or newer;
  \item optional: \code{mpi4py} 3.1 or newer, for the parallel $k$-point code
        (an MPI library such as OpenMPI must be installed first);
  \item optional, for development: \code{pytest}, \code{ruff} and
        \code{pre-commit}.
\end{itemize}

The dependencies are declared in \code{pyproject.toml}. \code{pip} installs
NumPy, SciPy and \pyscf{} automatically if they are missing.

\section{Installing from the repository}

It is recommended to install \pycmf{} into its own virtual environment:

\begin{lstlisting}[style=shell]
$ git clone https://github.com/QuantumLab-HCMUS/pyCMF-v1.git
$ cd pyCMF-v1
$ python3 -m venv .venv
$ source .venv/bin/activate
$ pip install --upgrade pip
$ pip install .                # users
$ pip install -e ".[test]"     # editable install with pytest
$ pip install -e ".[dev]"      # also ruff and pre-commit
$ pip install -e ".[mpi]"      # also mpi4py
\end{lstlisting}

With an editable install (\code{-e}) changes to the source code in
\code{src/pycmf} take effect without reinstalling.

\subsection*{Another PySCF on PYTHONPATH}

If the environment variable \code{PYTHONPATH} points to another checkout of
\pyscf{} (for example an older fork), that checkout takes precedence over the
\pyscf{} installed in the virtual environment. \pycmf{} then fails or, worse,
runs with the wrong version. Unset the variable before working with
\pycmf{}:

\begin{lstlisting}[style=shell]
$ unset PYTHONPATH
$ source .venv/bin/activate
$ python -c "import pyscf; print(pyscf.__version__, pyscf.__file__)"
\end{lstlisting}

The last command must print a version of at least 2.10 and a path inside
\code{.venv}. A single command can also be run with the variable removed,
\code{env -u PYTHONPATH .venv/bin/python script.py}.

\section{Checking the installation}

\begin{lstlisting}[style=shell]
$ python -c "import pycmf; print(pycmf.__version__)"
1.0.0
$ pytest -m "not slow"      # fast tests
$ pytest                    # all tests, including the slow k-point tests
\end{lstlisting}

The fast part of the suite has 106 tests and runs in a few minutes; the full
suite adds nine slower $k$-point tests (symmetry, larger meshes, MPI). Without
\code{mpi4py} the MPI test is skipped. A complete run on a 20-core
workstation took about 7~minutes for the fast tests and 2~minutes for the slow
ones.

The examples can be compared with their reference output:

\begin{lstlisting}[style=shell]
$ cd examples
$ python check_examples.py          # stable examples
$ python check_examples.py --all    # also the experimental ones
\end{lstlisting}

\section{Parallel execution}

\pycmf{} uses two levels of parallelism.

\begin{description}
  \item[Threads.] All large tensor contractions are matrix products that call
        the BLAS library of NumPy and the multithreaded routines of \pyscf{}.
        The number of threads is set with \code{OMP_NUM_THREADS}.
  \item[MPI.] The periodic code (\code{KOBMP2}) distributes the loop over
        $k$-points with \code{mpi4py}. The same script runs on every
        process; see Section~\ref{sec:kpt-mpi}.
\end{description}

The product of MPI processes and threads per process must not exceed the
number of physical cores. Oversubscribed BLAS threads slow a calculation down
by orders of magnitude. MPI can be switched off without uninstalling
\code{mpi4py} by setting the environment variable \code{PYCMF_NO_MPI=1}.

\section{Memory and scratch space}

The memory limit of a calculation is the attribute \code{max_memory} (in MB),
which is taken from the mean-field object. Large intermediate arrays are
written to the scratch directory of \pyscf{} (\code{TMPDIR} or
\code{pyscf.lib.param.TMPDIR}). Periodic calculations with density fitting
can need several gigabytes of disk space for the three-index integrals; see
Section~\ref{sec:kpt-resources}.

\section{Quick start}

The following script computes the OBMP2 energy of the OH radical and shows
the main stable methods.

\begin{lstlisting}[style=python]
from pyscf import gto, scf
import pycmf

mol = gto.M(atom='O 0 0 0; H 0 0 0.97', basis='cc-pvdz', spin=1)
mf = scf.UHF(mol).run()

mp = pycmf.OBMP2(mf).run()            # restricted or unrestricted, from mf
print(mp.e_tot, mp.spin_square())

o2 = pycmf.O2BMP2(mf, c_os=1.2).run() # c_os has no default

mf_df = scf.UHF(mol).density_fit().run()
dh = pycmf.OBDH(mf_df).run()          # B2PLYP by default

ham = pycmf.OBDF(mp, ncas=5, nelecas=(3, 2))   # active-space Hamiltonian
e_tot = ham.kernel()                           # FCI in the active space
h1, e_core = ham.get_h1eff()
h2 = ham.get_h2eff()
\end{lstlisting}

\chapter{Theory}
\label{ch:theory}

This chapter summarises what \pycmf{} computes, in the notation of the code.
The derivations are given in the papers listed in
Appendix~\ref{ch:references}.

\section{OBMP2}
\label{sec:theory-obmp2}

\subsection{Correlated Fock operator}

The electronic Hamiltonian $\hat H$ is transformed with an anti-Hermitian
doubles operator $\hat A$ built from first-order amplitudes. The
Baker--Campbell--Hausdorff (BCH) expansion is truncated at second order, with
the Fock operator $\hat F$ in the double commutator, and every commutator is
reduced to one-body operators and constants by a cumulant approximation,
denoted $[\ldots]_1$:
\begin{equation}
  \hat H_{\mathrm{OBMP2}} = \hat H_{\mathrm{HF}} + \big[\hat H,\hat A\big]_1
  + \tfrac12\big[[\hat F,\hat A],\hat A\big]_1
  = \hat{\bar F} + \bar C ,
  \qquad \bar f_{pq} = f_{pq} + v_{pq}.
  \label{eq:hobmp2}
\end{equation}
Here $f$ is the Hartree--Fock (HF) Fock matrix, $v$ is the
\emph{correlation potential} and $\fbar$ the \emph{correlated Fock matrix}.
The orbitals are the eigenvectors of $\fbar$ and the orbital energies
$\varepsilon_p$ its eigenvalues.

\subsection{Amplitudes and self-consistency}

The amplitudes are rebuilt from the current orbitals and orbital energies,
\begin{equation}
  t^{ab}_{ij} = c\,
  \frac{(ia|jb)}{\varepsilon_i+\varepsilon_j-\varepsilon_a-\varepsilon_b-\Delta},
  \label{eq:amplitudes}
\end{equation}
where $c$ is a scaling factor ($c=1$ in OBMP2) and $\Delta$ an optional level
shift (attribute \code{level_shift}). The cycle
\[
  \text{orbitals} \;\to\; t \;\to\; v \;\to\; \fbar \;\to\; \text{new orbitals}
\]
is repeated until the energy is stable and the Brillouin-type condition
$\bar f_{ia}=0$ holds. Starting from HF orbitals, the energy of the first
cycle is the conventional MP2 energy.

\subsection{Energy}

The energy of a cycle is
\begin{equation}
  E = E_{\mathrm{HF}}[\Phi] + c_0 + \sum_{i}^{\mathrm{occ}} v_{ii},
  \label{eq:energy}
\end{equation}
where $E_{\mathrm{HF}}[\Phi]$ is the HF energy of the current determinant and
$c_0$ the constant of the correlation potential. Because the amplitudes
depend on the orbital energies, this functional is not stationary with respect
to orbital rotations: the converged orbitals are defined by the eigenvalue
problem, not by a variational principle. The error of the energy is therefore
of first order in the residual $\bar f_{ia}$, and a tight convergence
threshold on $\bar f_{ia}$ (attribute \code{conv_tol_grad}) is needed when
energy differences below $10^{-6}$~Hartree matter.

\subsection{Spin blocks}

Amplitudes are stored in chemists' order, \code{t[i,a,j,b]}, as a list of
spin blocks $(x,y)$: the pair $(i,a)$ has spin $x$ and the pair $(j,b)$ has
spin $y$. The antisymmetrised amplitude is $\bar t = t - t^{\mathrm{exch}}$
for same-spin blocks and $\bar t = t$ for opposite-spin blocks. An
unrestricted reference has the four blocks $\alpha\alpha$, $\beta\beta$,
$\alpha\beta$ and $\beta\alpha$. A restricted reference has one spin channel
and two blocks (same spin, opposite spin); the spin-flipped partners
contribute a factor of two to the constants. The same routines
(\code{pycmf.mol.bch}) therefore serve both cases.

\subsection{Working equations}

With $f$ the HF Fock matrix in the current orbital basis and summation over
repeated indices, each block $(x,y)$ adds to the unsymmetrised matrices
$c_1^{x}$, $c_1^{y}$ and to the constant $c_0$.

\paragraph{First order, $[\hat H,\hat A]_1$.}
\begin{align}
  c_1^{y}[p,k] &\mathrel{+}= 2\,(ia|pl)\,\bar t_{ia,kl}, &
  c_1^{y}[p,l] &\mathrel{-}= 2\,(ia|kp)\,\bar t_{ia,kl}, \\
  c_1^{y}[j,b] &\mathrel{+}= 2\,f^{x}_{ia}\,\bar t_{ia,jb}, &
  c_0 &\mathrel{-}= (ia|jb)\,\bar t_{ia,jb}.
\end{align}

\paragraph{Second order, $\tfrac12[[\hat F,\hat A],\hat A]_1$.}
With the intermediates $V_{ia,jd} = f^{x}_{ac}\,\bar t_{ic,jd}$ and
$O_{ia,jd} = f^{x}_{ik}\,\bar t_{ka,jd}$,
\begin{align}
  c_1^{y}[j,k] &\mathrel{+}= t_{ia,jb}\,(V-O)_{ia,kb}, &
  c_1^{x}[i,k] &\mathrel{+}= t_{ia,jb}\,V_{ka,jb}, \\
  c_1^{x}[a,c] &\mathrel{+}= t_{ia,jb}\,O_{ic,jb}, &
  c_1^{y}[b,d] &\mathrel{+}= t_{ia,jb}\,(O-V)_{ia,jd}, \\
  c_0 &\mathrel{+}= t_{ia,jb}\,(O-V)_{ia,jb}, &
  c_1^{x}[i,a] &\mathrel{+}= \bar t_{ia,jb}\,y^{y}_{jb},
\end{align}
with $y^{y}_{jb}=\sum_{x'} f^{x'}_{kc}\,\bar t^{\,x'y}_{kc,jb}$, and, with
$R^{\mathrm{occ}}_{ik} = t_{ia,jb}\bar t_{ka,jb}$ and
$R^{\mathrm{vir}}_{ac} = t_{ia,jb}\bar t_{ic,jb}$ accumulated over the blocks
with first spin $x$,
\begin{equation}
  c_1^{x}[p,i] \mathrel{-}= f^{x}_{pk}R^{\mathrm{occ}}_{ik},\qquad
  c_1^{x}[p,c] \mathrel{-}= f^{x}_{pa}R^{\mathrm{vir}}_{ac}.
\end{equation}
Finally $v = \tfrac12\,(c_1 + c_1^{\mathsf T})$. The second-order part needs
no two-electron integrals.

\paragraph{First-order variant.}
\code{bch_order=1} keeps only the first commutator with half weight,
$\hat H + \tfrac12[\hat H,\hat A]_1$. It was the default of some modules of
version~0.2; the second-order form (\code{bch_order=2}) is the default now.

\subsection{Cost}

One cycle needs the integrals $(ia|jb)$, $(ia|pb)$ and $(ia|jp)$ and
contractions that scale as $o^2v^2N$ (first order) and $o^2v^3$ (second
order), where $o$, $v$ and $N$ are the numbers of occupied, virtual and all
orbitals. The cost is therefore $O(N^5)$ per cycle, like MP2. With exact
integrals the first index pair is transformed once per cycle,
$(ia|\mu\nu)$; with density fitting the three-index tensors $B^Q_{pa}$ and
$B^Q_{ip}$ are used. Memory is dominated by the amplitudes, $O(o^2v^2)$, and
by the three-index or half-transformed integrals.

\section{O2BMP2}
\label{sec:theory-o2bmp2}

Same-spin amplitudes are dropped and opposite-spin amplitudes are scaled,
$c_{\mathrm{SS}}=0$ and $c_{\mathrm{OS}}=c$. The scaling acts on the amplitudes
that enter $v$, so it changes the orbitals and orbital energies and not only
the energy. More generally, the attributes \code{c_os} and \code{c_ss} scale
the opposite-spin and same-spin amplitudes independently; $(1,1)$ is OBMP2,
$(0,0)$ is Hartree--Fock.

There is no universal value of $c_{\mathrm{OS}}$. The published choices are
1.2 (singlet--triplet gaps, charge-transfer excitations), 1.1 (ionization
potentials) and 1.7 (inverted singlet--triplet gaps). \code{pycmf.O2BMP2}
therefore has no default. An $O(N^4)$ algorithm based on the Laplace
transformation is not implemented.

\section{One-body double hybrids: OBDH and O2BDH}
\label{sec:theory-obdh}

For a double-hybrid functional with exact-exchange fraction $a_x$ and
correlation weight $a_c$,
\begin{equation}
  \fbar = f^{\mathrm{KS}}\big[a_x\,\mathrm{HF} + (1-a_x)\,X,\ (1-a_c)\,C\big]
          + a_c\,v_{\mathrm{OBMP2}},
  \qquad
  E = E^{\mathrm{KS}} + a_c\big(c_0 + \textstyle\sum_i v_{ii}\big),
\end{equation}
where $X$ and $C$ are the semilocal exchange and correlation functionals.
$v_{\mathrm{OBMP2}}$ is built with the HF Fock matrix of the current orbitals
and with amplitudes whose denominators are the eigenvalues of $\fbar$. With
$a_x=a_c=1$ the method is OBMP2; with $a_c=0$ it is the hybrid functional.
O2BDH uses opposite-spin amplitudes only, scaled by $c_{\mathrm{OS}}$.

The Kohn--Sham part is evaluated by \pyscf{} (libxc). The default functional
is B2PLYP. The global double hybrids of Table~\ref{tab:dh} are available by
name. Their parameters were fitted with a perturbative second-order energy on
Kohn--Sham orbitals and have not been re-optimised for the self-consistent
scheme.

\begin{table}[!htb]
\centering
\small
\caption{Double hybrids available by name (\code{pycmf.mol.obdh.FUNCTIONALS}).}
\label{tab:dh}
\begin{tabular}{@{}l l l l l l@{}}
\toprule
Name & $a_x$ & $a_c$ & Exchange & Correlation & Reference \\
\midrule
B2PLYP    & 0.53 & 0.27  & B88   & LYP & Grimme (2006) \\
mPW2PLYP  & 0.55 & 0.25  & mPW91 & LYP & Schwabe and Grimme (2006) \\
B2T-PLYP  & 0.60 & 0.31  & B88   & LYP & Tarnopolsky \emph{et al.} (2008) \\
B2GP-PLYP & 0.65 & 0.36  & B88   & LYP & Karton \emph{et al.} (2008) \\
B2K-PLYP  & 0.72 & 0.42  & B88   & LYP & Tarnopolsky \emph{et al.} (2008) \\
PBE0-DH   & 0.5  & 0.125 & PBE   & PBE & Br\'emond and Adamo (2011) \\
PBE-QIDH  & $3^{-1/3}$ & $1/3$ & PBE & PBE & Br\'emond \emph{et al.} (2014) \\
PBE0-2    & $2^{-1/3}$ & $1/2$ & PBE & PBE & Chai and Mao (2012) \\
\bottomrule
\end{tabular}
\end{table}

OBDH and O2BDH have not yet been described in a publication; this manual and
the docstring of \code{pycmf.mol.obdh} are their definition.

\section{Orbital-optimized double hybrids (OO-DH)}
\label{sec:theory-oodh}

OO-DH minimizes the double-hybrid energy itself with respect to the orbitals
(Peverati and Head-Gordon; OO-RI-B2PLYP of Neese \emph{et al.}):
\begin{equation}
  E[C] = E^{\mathrm{KS}}\big[a_x\,\mathrm{HF} + (1-a_x)\,X,\ (1-a_c)\,C;\ D\big]
         + a_c\,E_2[C],
\end{equation}
where $E_2$ is the MP2 energy built with the Kohn--Sham Fock matrix $f$ of the
same orbitals in semicanonical form (the occupied--occupied and
virtual--virtual blocks of $f$ are diagonalized, the singles are left out).
This is the OO-MP2 Lagrangian of Bozkaya \emph{et al.} with a Kohn--Sham
reference. In spin orbitals the Hylleraas functional is
\begin{equation}
  E_2 = \tfrac12\sum_{ijab} t_{ij}^{ab}\langle ij\|ab\rangle
        + \sum_{pq} f_{pq}\gamma_{pq},\qquad
  \gamma_{ij} = -\tfrac12\sum_{kab}t_{ik}^{ab}t_{jk}^{ab},\quad
  \gamma_{ab} = \tfrac12\sum_{ijc}t_{ij}^{ac}t_{ij}^{bc},
\end{equation}
and it is stationary in $t$. For the rotation
$\phi_i\to\phi_i+x_{ai}\phi_a$ the exact orbital gradient is therefore
\begin{equation}
  \frac{\partial E}{\partial x_{ai}} = 2f_{ai} + a_c\Big[
    \sum_{jbc}t_{ij}^{bc}\langle aj\|bc\rangle
  - \sum_{jkb}t_{jk}^{ab}\langle jk\|ib\rangle
  + 2\,(f\gamma-\gamma f)_{ai} + 2\,R[\gamma]_{ai}\Big],
\end{equation}
with $R[\gamma]$ the response of the Kohn--Sham Fock matrix to the density
$\gamma$ (Coulomb, scaled exact exchange and the exchange--correlation
kernel). The last term is what makes the gradient exact for a functional of
the density. The self-consistent driver diagonalizes
\begin{equation}
  \fbar = f + a_c v,\qquad
  v_{ai} = \frac{1}{2a_c}\Big(\frac{\partial E}{\partial x_{ai}} - 2f_{ai}\Big),
  \qquad v_{ij}=v_{ab}=0,
\end{equation}
so that the converged $\bar f_{ai}$ is half the orbital gradient. Because the
energy is stationary in the orbitals, $D + a_c\gamma$ is the relaxed density:
its expectation value of a one-electron operator is the derivative of the
energy (no $Z$-vector equations).

OO-DH differs from OBDH (Section~\ref{sec:theory-obdh}): OBDH adds the OBMP2
correlation potential to the Kohn--Sham Fock matrix and its energy is not
variational; OO-DH uses the exact gradient of the double-hybrid energy.
With $a_x=a_c=1$ and no semilocal parts it is OO-MP2; with $a_c=0$ it is the
hybrid functional. $c_{\mathrm{OS}}$ and $c_{\mathrm{SS}}$ scale the two spin
components of $E_2$.

\section{State-specific excited states}
\label{sec:theory-mom}

The correlated Fock operator is a one-body operator, so non-Aufbau
determinants can be optimised as in $\Delta$SCF. The calculation starts from a
MOM-HF determinant. After each diagonalisation the orbitals with the largest
overlap with the occupied space of the previous cycle (\code{'mom'}, Gilbert,
Besley and Gill) or of the initial determinant (\code{'imom'}, Barca, Gilbert
and Gill) are occupied. Excitation energies are differences of total
energies. An open-shell singlet obtained this way is a spin-mixed determinant
($\langle S^2\rangle\approx 1$).

\section{Transition moments and NOCI}
\label{sec:theory-noci}

The ground and excited determinants of a state-specific calculation have
different, non-orthogonal orbitals. For each spin the generalized
Slater--Condon rules are used in the corresponding-orbital form,
\begin{equation}
  X = C_L^{\mathsf T} S\, C_R = U\,\mathrm{diag}(\sigma)\,V^{\mathsf T},\qquad
  s = \det X,\qquad
  P = \pm\sum_i \Big[\prod_{k\neq i}\sigma_k\Big]\,|r_i\rangle\langle l_i|,
\end{equation}
with $l=C_LU$ and $r=C_RV$. The cofactor weights avoid a division by the
overlap, so the expressions also hold when $X$ is singular. The two-electron
terms use products of the same cofactors (Lemke, Kussmann and Ochsenfeld).

\paragraph{Transition dipole.} The two determinants are first made orthogonal
by the symmetric (L\"owdin) transformation (Bourne Worster, Feighan and
Manby), then $\boldsymbol\mu$ is taken between them and
$f = \tfrac23\,\Delta E\,|\boldsymbol\mu|^2$. This is a determinant-level
moment: the correlation enters through the orbitals only. For a closed-shell
molecule the excited determinant is spin mixed; \code{singlet=True}
multiplies $\boldsymbol\mu$ by $\sqrt2$, assuming the $\alpha$ and $\beta$
branches are equivalent and orthogonal.

\paragraph{NOCI.} The ground-state determinant and state-specific excited
determinants (one per excitation channel and spin) are combined by solving
$HC = SCE$ after canonical orthogonalization of $S$. Three Hamiltonians are
available:
\begin{itemize}
  \item \code{'bare'}: the physical Hamiltonian; the diagonal is the HF energy
        of each OBMP2 determinant.
  \item \code{'overlap_symmetric'}:
        $H_{IJ} = H^{\mathrm{bare}}_{IJ} + \tfrac12(\Delta_I+\Delta_J)S_{IJ}$
        with $\Delta_I = E_I - H^{\mathrm{bare}}_{II}$, so the diagonal is the
        OBMP2 energy.
  \item \code{'state_specific_fock'} (default): each determinant carries its
        correlated Fock operator $F_I$ and a constant $c_I$ with
        $c_I+\sum_{\mathrm{occ}}F_I=E_I$;
        $H_{IJ}=\tfrac12(c_I+c_J)S_{IJ}+\sum_\sigma\mathrm{tr}
        [\tfrac12(F_{I\sigma}+F_{J\sigma})D_{IJ\sigma}]$.
\end{itemize}
The last two are model choices; only their diagonal is fixed by the theory.
The averaged Fock operator is a one-body operator, so it does not couple the
$\alpha$ and $\beta$ determinants of one channel, which differ by two spin
orbitals (known issue~8).

\section{One-body downfolding (OBDF)}
\label{sec:theory-obdf}

Let $T_{\mathrm{act}}$ be the amplitudes with all four indices in an active
space. The external correlation potential and constant are
\begin{equation}
  v^{\mathrm{ext}} = v[T] - v[T_{\mathrm{act}}],\qquad
  c_0^{\mathrm{ext}} = c_0[T]-c_0[T_{\mathrm{act}}],
\end{equation}
and the active-space Hamiltonian is
\begin{equation}
  \hat H_{\mathrm{act}} = E_{\mathrm{core}}
  + \sum_{tu}\bar h_{tu}\,\hat a^\dagger_t\hat a_u
  + \tfrac12\sum_{tuvw}(tu|vw)\,\hat a^\dagger_t\hat a^\dagger_v\hat a_w\hat a_u,
\end{equation}
\begin{equation}
  \bar h = h + v^{\mathrm{HF}}_{\mathrm{core}} + v^{\mathrm{ext}},\qquad
  E_{\mathrm{core}} = E_{\mathrm{nuc}} + E^{\mathrm{HF}}_{\mathrm{core}}
  + c_0^{\mathrm{ext}} + \sum_{i\in\mathrm{core}} v^{\mathrm{ext}}_{ii}.
\end{equation}
Only the one-body integrals are dressed; the two-electron integrals are the
bare ones, so $\hat H_{\mathrm{act}}$ has the structure of a CASCI
Hamiltonian and can be handed to any solver (FCI, DMRG, selected CI, quantum
algorithms).

Two limits hold exactly and are tested: if the active space contains all
orbitals, $\hat H_{\mathrm{act}}=\hat H$ (FCI energy); if it is empty,
$E_{\mathrm{core}}=E_{\mathrm{OBMP2}}$. For any active space the reference
determinant has the energy
\begin{equation}
  \langle\Phi|\hat H_{\mathrm{act}}|\Phi\rangle
  = E_{\mathrm{OBMP2}}-E_c[T_{\mathrm{act}}],\qquad
  E_c[T_{\mathrm{act}}]=c_0[T_{\mathrm{act}}]+\sum_i v[T_{\mathrm{act}}]_{ii}:
\end{equation}
the correlation energy of the active double excitations is what is left to
the active-space solver.

\section{Properties}
\label{sec:theory-properties}

\begin{itemize}
  \item \emph{HF-like density}: the density of the OBMP2 determinant.
  \item \emph{MP2-like density}: $\gamma_{\mathrm{oo}}=1-R^{\mathrm{occ}}$ and
        $\gamma_{\mathrm{vv}}=R^{\mathrm{vir}}$ with the converged amplitudes
        (unrelaxed). After one cycle from HF orbitals it equals the unrelaxed
        MP2 density.
  \item $\langle S^2\rangle$ of the determinant, optionally with the
        first-order doubles correction
        $\langle S^2\rangle = \langle S^2\rangle_{\mathrm{det}}
        - \sum_{ia,jb} S_{ib}S_{ja}t_{ia,jb}$, where $S$ is the overlap of
        the $\alpha$ and $\beta$ orbitals.
  \item Dipole moments and isotropic hyperfine coupling constants as
        expectation values over either density.
  \item Koopmans-type ionization potentials and electron affinities from the
        eigenvalues of $\fbar$.
\end{itemize}

\section{Projection embedding in DFT}
\label{sec:theory-embed}

In the projection-based embedding (\code{pycmf.embed}, experimental) the two
levels have their own functionals. A Kohn--Sham calculation of the full
system with the low-level functional \code{xc_env} gives the densities
$\gamma_A$ and $\gamma_B$ of the subsystem and of the environment (SPADE
partition) and the embedding potential
\begin{equation}
  v_{\mathrm{emb}} = v^{\mathrm{env}}[\gamma_A+\gamma_B]-v^{\mathrm{env}}[\gamma_A]
  +\mu\,S\gamma_BS .
\end{equation}
The subsystem is then treated with OBMP2 or OBDH (with its own double hybrid)
in this potential, optionally in a virtual space truncated by concentric
localization, and the total energy is
\begin{equation}
  E = E^{\mathrm{high}}[\gamma_A'] + E^{\mathrm{env}}[\gamma_A'+\gamma_B]
  -E^{\mathrm{env}}[\gamma_A'] + \mu\,\mathrm{tr}(\gamma_A'S\gamma_BS),
\end{equation}
with $\gamma_A'$ the relaxed density of the embedded subsystem. Two limits are
exact and tested: with all atoms in the subsystem $E$ is the energy of the
high-level method, and with the same functional on both levels and $a_c=0$ it
is the Kohn--Sham energy of the full system.

\section{Periodic systems}
\label{sec:theory-pbc}

For a crystal with a $k$-point mesh of $\Nk$ points the working equations are
those of Section~\ref{sec:theory-obmp2}, spin-summed and with conservation
of crystal momentum. For the four $k$-points $(k_i,k_a,k_j,k_b)$ of a double
excitation $ia\to jb$,
\begin{equation}
  g_{iajb} = \frac{(ia|jb)}{\Nk},\qquad
  t_{iajb} = \frac{g^*_{iajb}}{\varepsilon_i-\varepsilon_a+\varepsilon_j-\varepsilon_b-\Delta},\qquad
  \bar t_{iajb} = \alpha\,t_{iajb} - \beta\,t_{ibja},
\end{equation}
with $\alpha = (c_{\mathrm{SS}}+c_{\mathrm{OS}})/2$ and $\beta =
c_{\mathrm{SS}}/2$ in the terms linear in the amplitudes, and the squares of
the scaling factors in the quadratic terms. All two-electron integrals are
products of Gaussian density-fitting tensors,
\begin{equation}
  (p k_1, q k_2 \,|\, r k_3, s k_4) = \sum_Q L^Q_{pq}(k_1,k_2)\,L^Q_{rs}(k_3,k_4),
\end{equation}
so no four-index transformation is performed. The algorithm, the use of
$k$-point symmetry and the quasiparticle methods built on KOBMP2 are
described in Chapter~\ref{ch:periodic}.

\chapter{Molecular Calculations}
\label{ch:molecular}

All molecular methods follow the same pattern: build the object from a
converged mean-field object of \pyscf{}, set attributes, call
\code{kernel()} or \code{run()}, and read the results from attributes.

\section{Choosing the method}

\begin{table}[!htb]
\centering
\small
\caption{Factory functions for molecules.}
\label{tab:factories}
\begin{tabularx}{\textwidth}{@{}l X l@{}}
\toprule
Call & Reference and class & Result \\
\midrule
\code{pycmf.OBMP2(mf)} & RHF or RKS $\to$ \code{ROBMP2}; UHF or ROHF $\to$ \code{UOBMP2} & OBMP2 \\
\code{pycmf.O2BMP2(mf, c_os=1.2)} & as above; \code{c_os} is mandatory & O2BMP2 \\
\code{pycmf.OBDH(mf, xc='B2PLYP')} & as above $\to$ \code{ROBDH} or \code{UOBDH} & one-body double hybrid \\
\code{pycmf.O2BDH(mf, c_os=1.2)} & as above; \code{c_os} is mandatory & opposite-spin OBDH \\
\code{pycmf.OBDF(mp, ncas, nelecas)} & converged OBMP2 or O2BMP2 object & active-space Hamiltonian \\
\bottomrule
\end{tabularx}
\end{table}

The factory functions choose the restricted or unrestricted class from the
type of the mean-field object. A ROHF reference is converted to UHF form.
Keyword arguments of the factory functions are set as attributes, so
\begin{lstlisting}[style=python]
mp = pycmf.OBMP2(mf, conv_tol=1e-10, level_shift=0.1)
\end{lstlisting}
is the same as setting \code{mp.conv_tol} and \code{mp.level_shift} after
construction. The classes \code{ROBMP2}, \code{UOBMP2}, \code{ROBDH} and
\code{UOBDH} can also be instantiated directly.

\section{OBMP2 and O2BMP2}
\label{sec:obmp2}

\subsection{A first calculation}

\begin{lstlisting}[style=python]
from pyscf import gto, scf, mp
import pycmf

mol = gto.M(atom='O 0 0 0.1173; H 0 0.7572 -0.4692; H 0 -0.7572 -0.4692',
            basis='cc-pvdz')
mf = scf.RHF(mol).run()
ob = pycmf.OBMP2(mf).run()
print(ob.e_tot, ob.converged)
print(mp.MP2(mf).run().e_tot)          # conventional MP2 for comparison
print(ob.ip_koopmans() * 27.2114)       # Koopmans IP in eV
\end{lstlisting}

For water in the cc-pVDZ basis this gives (from
\code{examples/01-obmp2.py}):

\begin{lstlisting}[style=output]
E(RHF)    = -76.0267720534
E(MP2)    = -76.2307756326
E(OBMP2)  = -76.2312896416   converged: True
Koopmans IP (eV) = 14.953
\end{lstlisting}

\subsection{Attributes}

\begin{longtable}{@{}p{0.27\textwidth} p{0.16\textwidth} p{0.49\textwidth}@{}}
\caption{Attributes of the molecular OBMP2 classes.}\label{tab:attrs}\\
\toprule
Attribute & Default & Meaning \\
\midrule
\endfirsthead
\toprule
Attribute & Default & Meaning \\
\midrule
\endhead
\bottomrule
\endfoot
\code{conv_tol} & \code{1e-8} & convergence threshold of the energy change between cycles \\
\code{conv_tol_grad} & $\sqrt{\text{\code{conv_tol}}}$ & threshold on the norm of the occupied--virtual block $\bar f_{ia}$ \\
\code{max_cycle} & \code{100} & maximum number of cycles; \code{max_cycle=1} from HF orbitals gives conventional MP2 \\
\code{bch_order} & \code{2} & 2: second-order BCH; 1: half-weighted first commutator \\
\code{c_os}, \code{c_ss} & \code{1, 1} & scaling of the opposite-spin and same-spin amplitudes \\
\code{level_shift} & \code{0} & constant added to the energy gap in the amplitude denominators \\
\code{occupation} & \code{'aufbau'} & \code{'mom'} or \code{'imom'} for state-specific calculations \\
\code{diis} & \code{True} & DIIS extrapolation of the correlated Fock matrix \\
\code{diis_space} & \code{8} & number of stored DIIS vectors \\
\code{diis_start_cycle} & \code{1} & first cycle with DIIS \\
\code{with_df} & from \code{mf} & density-fitting object; \code{None} means exact integrals \\
\code{max_memory} & from \code{mf} & memory limit in MB \\
\code{verbose}, \code{stdout} & from \code{mol} & output level and stream (\pyscf{} logger) \\
\end{longtable}

\code{frozen} orbitals are not supported for molecules: all orbitals are
relaxed by the correlated Fock operator, and a value other than \code{None}
or \code{0} raises \code{NotImplementedError}.

\subsection{Results}

After \code{kernel()} the following attributes are set.

\begin{description}[style=nextline,leftmargin=3.2cm,labelwidth=3cm]
  \item[\code{converged}] \code{True} if both convergence criteria were met.
  \item[\code{e_tot}] total energy; also the return value of \code{kernel()}.
  \item[\code{e_corr}] \code{e_tot} minus the energy of the input mean field.
  \item[\code{e_ref}] HF (or, for OBDH, Kohn--Sham) energy of the final determinant.
  \item[\code{e_hf}] energy of the input mean field.
  \item[\code{mo_energy}] eigenvalues of the correlated Fock matrix.
  \item[\code{mo_coeff}] correlated orbitals, occupied orbitals first.
  \item[\code{mo_occ}] occupation numbers in the same order.
\end{description}

For an unrestricted reference these arrays have a leading spin index of
length two, as in \pyscf{}.

\subsection{Methods of the result object}

\begin{longtable}{@{}p{0.40\textwidth} p{0.54\textwidth}@{}}
\toprule
Method & Returns \\
\midrule
\endhead
\bottomrule
\endfoot
\code{get_fock(ao_repr=False)} & converged correlated Fock matrix (MO or AO basis) \\
\code{get_vcorr()} & correlation potential $v$ in the basis of the converged orbitals \\
\code{make_rdm1(ao_repr=True, correlated=False)} & one-particle density; \code{correlated=True} gives the MP2-like density \\
\code{dip_moment(correlated=False, unit='Debye')} & dipole moment over either density \\
\code{ip_koopmans()}, \code{ea_koopmans()} & $-\varepsilon_{\mathrm{HOMO}}$ and $-\varepsilon_{\mathrm{LUMO}}$ of $\fbar$ (per spin for unrestricted) \\
\code{spin_square(correlated=False)} & $\langle S^2\rangle$ and $2S+1$ (unrestricted only; \code{ROBMP2} returns \code{(0, 1)}) \\
\code{hfcc(atoms=None, correlated=False)} & isotropic hyperfine couplings in MHz (unrestricted only) \\
\code{get_intermediates(with_amplitudes=False)} & quantities of the last cycle (\code{Intermediates}); amplitudes are rebuilt on request \\
\code{density_fit(auxbasis=None, with_df=None)} & switches the correlation part to density fitting (returns \code{self}) \\
\end{longtable}

\subsection{O2BMP2 and spin-component scaling}

\begin{lstlisting}[style=python]
o2 = pycmf.O2BMP2(mf, c_os=1.2).run()            # c_ss = 0
scs = pycmf.OBMP2(mf, c_os=1.2, c_ss=1/3).run()  # any other scaling
\end{lstlisting}

\code{O2BMP2} raises \code{ValueError} if \code{c_os} is not given. With
density fitting an O2BMP2 cycle costs about 60--70\,\% of an OBMP2 cycle
because the same-spin blocks are skipped. The example
\code{examples/02-o2bmp2.py} shows the effect of \code{c_os} on the
singlet--triplet gap of CH$_2$ (cc-pVTZ, density fitting):

\begin{lstlisting}[style=output]
method             E(S)-E(T) kcal/mol
HF                      28.24
OBMP2                   16.16
O2BMP2 c_os=1.1         12.09
O2BMP2 c_os=1.2         12.35
O2BMP2 c_os=1.3         12.90
\end{lstlisting}

\subsection{Open shells}

An unrestricted reference gives unrestricted OBMP2. OBMP2 reduces the spin
contamination of UHF considerably. For the CN radical in the cc-pVDZ basis
(\code{examples/06-properties.py}), $\langle S^2\rangle$ drops from 1.150
(UHF) to 0.766 (OBMP2 determinant) and 0.761 with the amplitude correction.

\begin{lstlisting}[style=python]
mol = gto.M(atom='O 0 0 0; H 0 0 0.97', basis='cc-pvdz', spin=1)
mf = scf.UHF(mol).run()
ob = pycmf.OBMP2(mf).run()
print(ob.e_tot, ob.spin_square())
\end{lstlisting}

\section{Density fitting}
\label{sec:df}

There are two ways to use density-fitted integrals.

\begin{enumerate}
  \item The mean field uses density fitting:
        \code{mf = scf.RHF(mol).density_fit()}. \pycmf{} reuses
        \code{mf.with_df} (JK-fit basis) for the correlation part, as the
        DF-MP2 of \pyscf{} does.
  \item The mean field uses exact integrals:
        \code{pycmf.OBMP2(mf).density_fit()} fits only the correlation part,
        with an RI-C auxiliary basis chosen by \pyscf{};
        \code{density_fit(auxbasis='cc-pvtz-ri')} selects a specific one.
\end{enumerate}

\begin{lstlisting}[style=python]
mf_df = scf.RHF(mol).density_fit().run()
e1 = pycmf.OBMP2(mf_df).kernel()                              # JK-fit basis
e2 = pycmf.OBMP2(mf).density_fit(auxbasis='cc-pvdz-ri').kernel()   # RI-C basis
\end{lstlisting}

For water/cc-pVDZ the two give $-76.2312271301$ and $-76.2312761456$~Hartree,
against $-76.2312896416$ with exact integrals. Exact integrals need the AO
integrals in memory; for more than a few hundred basis functions use density
fitting.

\section{Convergence}
\label{sec:convergence}

Every method uses the same self-consistent driver
(\code{pycmf.core.scf_loop}). At each cycle it

\begin{enumerate}
  \item builds the correlated Fock matrix from the current orbitals and
        orbital energies;
  \item extrapolates it with DIIS, using the commutator $[F,D]$ per spin
        channel as error vector;
  \item diagonalises it in the space spanned by the current orbitals;
  \item occupies the new orbitals according to \code{occupation}.
\end{enumerate}

Convergence requires both $|\Delta E| <$ \code{conv_tol} and
$\|\bar f_{ia}\| <$ \code{conv_tol_grad}. At \code{verbose >= 4} the
settings are printed first, then the energy, $\Delta E$ and
$\|\bar f_{ia}\|$ of every cycle. For water/cc-pVDZ with the default settings:

\begin{lstlisting}[style=output]
******** <class 'pycmf.mol.obmp2.ROBMP2'> ********
bch_order = 2
c_os = 1  c_ss = 1
level_shift = 0
occupation = aufbau
conv_tol = 1e-08
max_cycle = 100
DIIS = True, space = 8
integrals = exact
cycle= 1 E= -76.2307756295253  delta_E= -76.2  |f_ov|= 0.0302
cycle= 2 E= -76.231238968644  delta_E= -0.000463  |f_ov|= 0.00485
cycle= 3 E= -76.2312803148426  delta_E= -4.13e-05  |f_ov|= 0.00109
cycle= 4 E= -76.2312879484375  delta_E= -7.63e-06  |f_ov|= 0.000263
cycle= 5 E= -76.2312891373978  delta_E= -1.19e-06  |f_ov|= 9.69e-05
cycle= 6 E= -76.231289619696  delta_E= -4.82e-07  |f_ov|= 1.47e-05
cycle= 7 E= -76.2312896408706  delta_E= -2.12e-08  |f_ov|= 1.6e-06
cycle= 8 E= -76.2312896416401  delta_E= -7.69e-10  |f_ov|= 1.93e-07
ROBMP2 converged.  E(ROBMP2) = -76.2312896416401
\end{lstlisting}

The first cycle reproduces the MP2 energy ($-76.23077563$) to the accuracy
of the HF reference.

If a calculation does not converge:

\begin{itemize}
  \item increase \code{max_cycle};
  \item try a small \code{level_shift} (for example 0.1--0.3~Hartree), which
        damps the amplitudes of near-degenerate pairs;
  \item start DIIS later (\code{diis_start_cycle}) or switch it off to see
        whether the iterations oscillate;
  \item check that the mean field is a stable solution
        (\code{mf.stability()} in \pyscf{}).
\end{itemize}

A calculation that stops at \code{max_cycle} prints a warning and sets
\code{converged=False}; its energy is that of the last cycle.

\section{One-body double hybrids}
\label{sec:obdh}

\code{OBDH} and \code{O2BDH} take a converged HF or Kohn--Sham object, which
provides the starting orbitals. The functional is defined in one of three
ways:

\begin{lstlisting}[style=python]
pycmf.OBDH(mf)                                           # B2PLYP
pycmf.OBDH(mf, xc='B2GP-PLYP')                           # by name
pycmf.OBDH(mf, ax=0.6, ac=0.3, xc_x='PBE', xc_c='PBE')   # four parameters
pycmf.OBDH(mf, xc='0.5*HF + 0.5*B88, 0.75*LYP', ac=0.25) # PySCF string + a_c
pycmf.O2BDH(mf, c_os=1.2, xc='PBE0-DH')
\end{lstlisting}

Names are case-insensitive, and hyphens and underscores are ignored. With a
functional string the semilocal correlation must already be scaled by
$1-a_c$, and \code{ax} must not be given (the HF fraction is part of the
string). Explicit \code{ax} and \code{ac} override the parameters of a named
functional.

\begin{longtable}{@{}p{0.22\textwidth} p{0.18\textwidth} p{0.52\textwidth}@{}}
\caption{Additional attributes of \code{ROBDH} and \code{UOBDH}.}\\
\toprule
Attribute & Default & Meaning \\
\midrule
\endhead
\bottomrule
\endfoot
\code{ax} & \code{0.53} & fraction of HF exchange \\
\code{ac} & \code{0.27} & weight of the OBMP2 correlation \\
\code{xc_x}, \code{xc_c} & \code{'B88'}, \code{'LYP'} & semilocal exchange and correlation \\
\code{xc} & \code{None} & name of a double hybrid or a \pyscf{} functional string \\
\code{grids_level} & \code{None} & DFT integration grid level (default of \pyscf{}) \\
\end{longtable}

Useful methods: \code{get_xc()} returns the functional string of the
Kohn--Sham part, \code{set_functional(name)} takes the four parameters from a
named double hybrid, and \code{get_ks()} returns the Kohn--Sham object that
evaluates the hybrid part. When the reference uses density fitting, the
Kohn--Sham part uses the same density-fitting object.

From \code{examples/05-obdh.py} (water, cc-pVDZ, density fitting):

\begin{lstlisting}[style=output]
functional          : 0.53*HF + 0.47*B88, 0.73*LYP  + 0.27 E_c(OBMP2)
E(OBDH)             = -76.3425040639   converged: True
E(O2BDH, c_os=1.2)  = -76.3231558818
B2GP-PLYP  a_x = 0.6500  a_c = 0.3600  E = -76.3273800422
PBE0-DH    a_x = 0.5000  a_c = 0.1250  E = -76.3282840566
PBE-QIDH   a_x = 0.6934  a_c = 0.3333  E = -76.3072423903
\end{lstlisting}

For degenerate open shells such as the $^2\Pi$ state of OH, the orientation
of the open-shell orbital is fixed only by the numerical noise of the DFT
grid. Convergence then stalls at $\|\bar f_{ia}\|\sim10^{-6}$ and total
energies are reproducible to about $10^{-7}$~Hartree only.

\section{State-specific excited states}
\label{sec:mom}

An excited state is computed from a non-Aufbau determinant. First converge
the excited determinant at the HF level with the MOM of \pyscf{}, then run
OBMP2 or O2BMP2 with \code{occupation='imom'}:

\begin{lstlisting}[style=python]
mol = gto.M(atom='O 0 0 0.1173; H 0 0.7572 -0.4692; H 0 -0.7572 -0.4692',
            basis='aug-cc-pvdz')
mf_gs = scf.UHF(mol).run()

occ = mf_gs.mo_occ.copy()                   # alpha HOMO -> LUMO
homo = mol.nelec[0] - 1
occ[0][homo], occ[0][homo + 1] = 0, 1
mf_ex = scf.addons.mom_occ(scf.UHF(mol), mf_gs.mo_coeff, occ)
mf_ex.kernel(mf_ex.make_rdm1(mf_gs.mo_coeff, occ))

gs = pycmf.O2BMP2(mf_gs, c_os=1.2).run()
ex = pycmf.O2BMP2(mf_ex, c_os=1.2, occupation='imom').run()
omega = (ex.e_tot - gs.e_tot) * 27.2114     # eV
\end{lstlisting}

The occupation numbers \code{mo_occ} of the mean-field object define the
starting determinant; the orbitals are reordered occupied-first internally.
At \code{verbose >= 5} the number of electrons that have left the reference
occupied space is printed every cycle, which reveals a collapse to another
state. \code{'imom'} keeps the reference fixed and is usually more robust
than \code{'mom'}. For the example above (\code{examples/03-excited_state.py}):

\begin{lstlisting}[style=output]
method             omega (eV)    <S^2>
MOM-HF                  6.209    1.008
OBMP2                   7.588    1.006   converged: True
O2BMP2 (1.2)            7.193    1.006   converged: True
\end{lstlisting}

The excited determinant is a 50:50 mixture of singlet and triplet. Spin
purification is left to the user.

\section{One-body downfolding}
\label{sec:obdf}

\code{OBDF} builds the active-space Hamiltonian of
Section~\ref{sec:theory-obdf} from a converged OBMP2 or O2BMP2 object (OBDH
objects are not accepted).

\begin{lstlisting}[style=python]
mp = pycmf.OBMP2(mf).run()
ham = pycmf.OBDF(mp, ncas=6, nelecas=6)        # orbitals around the Fermi level
ham = pycmf.OBDF(mp, 4, (1, 1), cas_list=[0, 1, 4, 5])   # 0-based indices
e = ham.kernel()                  # FCI in the active space
h1, e_core = ham.get_h1eff()      # dressed one-body integrals and constant
h2 = ham.get_h2eff()              # bare active-space integrals
ham.to_fcidump('active.fcidump')  # restricted reference only
\end{lstlisting}

\begin{description}[style=nextline,leftmargin=3.2cm,labelwidth=3cm]
  \item[\code{ncas}] number of active orbitals.
  \item[\code{nelecas}] number of active electrons, an integer or a pair
        $(n_\alpha, n_\beta)$.
  \item[\code{cas_list}] 0-based indices of the active orbitals in the orbital
        order of \code{mp}; a pair of lists for different $\alpha$ and
        $\beta$ choices. By default the \code{ncas} orbitals around the Fermi
        level that hold \code{nelecas} electrons are used.
  \item[\code{fcisolver}] solver with the interface of \code{pyscf.fci};
        it can be replaced by any compatible solver (DMRG, selected CI).
  \item[\code{nroots}] number of states requested from the solver.
\end{description}

Results: \code{e_core}, \code{h1eff}, \code{h2eff}, and after
\code{kernel()} the total energy \code{e_tot}, the active-space energy
\code{e_cas} and the CI vector(s) \code{ci}. For an unrestricted reference
\code{h1eff} is a pair and \code{h2eff} a triple ($\alpha\alpha$,
$\alpha\beta$, $\beta\beta$), and \code{to_fcidump} is not available.

Do not cut through a set of degenerate orbitals with \code{cas_list}: the
result then depends on an arbitrary rotation. \code{OBDF} warns in this
situation.

\code{examples/04-obdf.py} compares OBDF(6,6) with OBMP2, CASCI(6,6) and
CCSD(T) along the dissociation curve of N$_2$ (cc-pVDZ):

\begin{lstlisting}[style=output]
 R (A)          OBMP2     CASCI(6,6)      OBDF(6,6)        CCSD(T)
  1.00    -109.215012    -108.980201    -109.220021    -109.231138
  1.10    -109.267377    -109.021905    -109.272409    -109.279398
  1.30    -109.229093    -108.963955    -109.225534    -109.224738
  1.60    -109.122703    -108.828238    -109.073809    -109.076363
  2.00    -109.053390    -108.741134    -108.924800    -109.015388
\end{lstlisting}

\section{Properties}
\label{sec:properties}

\begin{lstlisting}[style=python]
mp.make_rdm1(ao_repr=True, correlated=False)  # HF-like; True: MP2-like
mp.dip_moment(correlated=True)
mp.spin_square(correlated=False)              # unrestricted only
mp.hfcc(atoms=None, correlated=False)         # isotropic couplings, MHz
mp.ip_koopmans(); mp.ea_koopmans()
mp.get_fock(); mp.get_vcorr()                 # MO basis
\end{lstlisting}

The functions behind these methods are in \code{pycmf.properties}
(\code{rdm}, \code{spin}, \code{hfcc}). From
\code{examples/06-properties.py} (CN, cc-pVDZ):

\begin{lstlisting}[style=output]
<S^2>: UHF 1.1497 | OBMP2 determinant 0.7657 | with amplitudes 0.7609
HF-like  density: dipole 1.2327 D   A_iso(13C)   547.33 MHz   A_iso(14N)   -18.43 MHz
MP2-like density: dipole 1.0926 D   A_iso(13C)   560.71 MHz   A_iso(14N)   -10.45 MHz
\end{lstlisting}

\section{Experimental: restricted open-shell OBMP2}
\label{sec:roobmp2}

\code{ROOBMP2} works with a high-spin ROHF reference and keeps common spatial
orbitals, ordered closed, open, virtual. It has no formal derivation yet, no
DIIS and no state following.

\begin{lstlisting}[style=python]
from pycmf.mol.roobmp2 import ROOBMP2
mf = scf.ROHF(mol).run()
ro = ROOBMP2(mf).run()
print(ro.e_tot, ro.spin_square())      # <S^2> = 0.75 for a doublet
\end{lstlisting}

The additional attribute \code{conv_tol_denom} (default $10^{-6}$) is the
threshold on the difference between the amplitude denominators and the
diagonals of the spin Fock matrices. \code{mo_energy} carries the spin
energies as the tags \code{mo_ea} and \code{mo_eb}.

\section{Experimental: embedding in DFT}
\label{sec:embedding}

\begin{lstlisting}[style=python]
from pyscf import dft
from pycmf.embed import EmbeddedOBDH, EmbeddedOBMP2

ks = dft.UKS(mol, xc='pbe').density_fit().run()   # low level, full system
emb = EmbeddedOBDH(ks, active_atoms=[0, 1, 2], xc='B2PLYP',
                   use_cl=True, n_shells=1)
e = emb.kernel()
emb.e_env, emb.e_sub, emb.nelec_sub
emb.make_rdm1(); emb.dip_moment()                 # full system
EmbeddedOBMP2(ks, active_atoms=[0, 1, 2]).kernel()
\end{lstlisting}

\begin{longtable}{@{}p{0.22\textwidth} p{0.18\textwidth} p{0.52\textwidth}@{}}
\caption{Attributes of \code{EmbeddedOBMP2} and \code{EmbeddedOBDH}.}\\
\toprule
Attribute & Default & Meaning \\
\midrule
\endhead
\bottomrule
\endfoot
\code{active_atoms} & \code{None} & indices of the atoms of the subsystem \\
\code{xc_env} & from \code{mf} & low-level functional (environment and embedding potential); \code{'pbe'} for an HF object \\
\code{mu} & \code{1e6} & level-shift parameter of the projector \\
\code{use_cl} & \code{False} & truncate the virtual space by concentric localization \\
\code{n_shells} & \code{1} & number of CL shells kept \\
\end{longtable}

Results: \code{e_tot} (embedded total energy), \code{e_env} (low-level
energy of the full system), \code{e_sub} (high-level energy of the
subsystem), \code{e_low_sub}, \code{nelec_sub}, and the orbitals of the
embedded subsystem. \code{make_rdm1(with_env=True)} returns the density of
the full system. The embedding is formulated for unrestricted references
only. The helper functions \code{spade_partition}, \code{atom_aos},
\code{concentric_localization} and \code{shell_sizes} are exported by
\code{pycmf.embed}.

In version 1.0 the two levels have separate functionals: the environment and
the embedding potential use \code{xc_env}, and the subsystem uses its own
double hybrid. In version 0.2 \code{xc_env} was also used for the
Kohn--Sham part of the subsystem; those numbers are reproduced with
\code{xc=xc_env}.

\chapter{Periodic Systems}
\label{ch:periodic}

The periodic code in \code{pycmf.pbc} is \textbf{experimental}. It
reproduces the previous implementation (version 0.2) and the KMP2 code of
\pyscf{} at the first cycle, but it has not yet been compared with published
numbers. Every class of the sub-package emits an \code{ExperimentalWarning}
when it is instantiated.

\begin{lstlisting}[style=python]
from pycmf.pbc import KOBMP2, KO2BMP2, KOBDF, KUOBMP2
from pycmf.pbc import remove_linear_dependency_
\end{lstlisting}

\section{k-point OBMP2: KOBMP2 and KO2BMP2}
\label{sec:kobmp2}

\subsection{Usage}

\begin{lstlisting}[style=python]
from pyscf.pbc import gto, scf
from pycmf.pbc import KOBMP2, KO2BMP2

a0 = 3.5668
cell = gto.M(atom='C 0 0 0; C %f %f %f' % (a0/4, a0/4, a0/4),
             a=[[0, a0/2, a0/2], [a0/2, 0, a0/2], [a0/2, a0/2, 0]],
             basis='gth-szv', pseudo='gth-pade')
mf = scf.KRHF(cell, cell.make_kpts([2, 2, 2])).density_fit().run(conv_tol=1e-10)

mp = KOBMP2(mf).run()            # e_tot, e_corr per unit cell
print(mp.e_tot, mp.band_gap())   # gap of the bands of the correlated Fock matrix

o2 = KO2BMP2(mf, c_os=1.2).run() # c_os has no default
KOBMP2(mf, kpt_symmetry=False)   # full k-mesh, no symmetry
\end{lstlisting}

With the diamond cell above (\code{examples/experimental/kobmp2.py}):

\begin{lstlisting}[style=output]
E(KRHF)    = -10.93208054   gap = 0.67421
E(KOBMP2)  = -11.02688135   gap = 0.73499
without k-point symmetry: 8 k-points instead of 3, same energy: True
E(KO2BMP2, c_os = 1.2) = -11.00812694   gap = 0.72361
\end{lstlisting}

\subsection{Requirements on the mean field}

\begin{itemize}
  \item The reference must be a closed-shell \code{KRHF} object with integer
        occupations. Kohn--Sham objects are rejected. The number of
        occupied bands may depend on the $k$-point.
  \item The system must be three-dimensional.
  \item The correlation part uses Gaussian density fitting. If the mean field
        uses a \code{GDF} or \code{RSGDF} object, it is reused; otherwise
        (for example FFT-based integrals) a \code{GDF} object is built for
        the correlation part, and the result differs from an all-FFT
        calculation by the density-fitting error.
  \item A mean field that uses $k$-point symmetry (\code{KPoints} objects of
        \pyscf{}) is converted to the full mesh.
\end{itemize}

\subsection{Attributes}

\begin{longtable}{@{}p{0.28\textwidth} p{0.17\textwidth} p{0.47\textwidth}@{}}
\caption{Attributes of \code{KOBMP2}.}\label{tab:kattrs}\\
\toprule
Attribute & Default & Meaning \\
\midrule
\endfirsthead
\toprule
Attribute & Default & Meaning \\
\midrule
\endhead
\bottomrule
\endfoot
\code{conv_tol} & \code{1e-8} & convergence of the energy per cell \\
\code{conv_tol_grad} & $\sqrt{\text{\code{conv_tol}}}$ & convergence of the occupied--virtual block of $\fbar$ \\
\code{max_cycle} & \code{100} & maximum number of cycles \\
\code{diis}, \code{diis_space}, \code{diis_start_cycle} & \code{True, 8, 1} & DIIS on the Fock matrices of the irreducible $k$-points \\
\code{bch_order} & \code{2} & 1: half-weighted first commutator \\
\code{c_os}, \code{c_ss} & \code{1, 1} & amplitude scaling; \code{KO2BMP2} sets \code{c_ss = 0} \\
\code{level_shift} & \code{0} & shift of the energy denominators \\
\code{frozen} & \code{None} & number of lowest bands, or list of band indices, without amplitudes \\
\code{kpt_symmetry} & \code{True} & \code{True}: space group and time reversal; \code{'time_reversal'}; \code{False} \\
\code{with_df} & from \code{mf} & GDF or RSGDF object; \code{mp.density_fit(auxbasis=...)} builds another one \\
\code{max_memory} & from \code{mf} & decides whether three-index integrals stay in memory \\
\code{direct_exchange} & \code{True} & take the HF exchange matrices from the pass over the integrals that builds the amplitudes \\
\code{single_precision} & \code{False} & store the occupied--virtual three-index tensors in single precision \\
\end{longtable}

\code{frozen} and \code{kpt_symmetry} can also be given to the constructor,
\code{KOBMP2(mf, frozen=None, kpt_symmetry=None)}. Frozen orbitals are
excluded from the amplitudes at every $k$-point but are still relaxed by the
correlated Fock operator. The frozen set must not cut a degenerate level; in
diamond the two lowest valence bands are degenerate at X, so \code{frozen=1}
is not valid there (a warning is issued).

\code{direct_exchange} gives identical energies and saves one pass over the
integral file per cycle; it applies to a plain \code{KRHF} with GDF and
without MPI. \code{single_precision} halves the memory of the stored tensors;
the arithmetic stays in double precision, and the energy changes by less than
$10^{-9}$~Hartree for diamond gth-dzvp $3\times3\times3$.

The names of version 0.2 are kept as aliases: \code{niter}
(\code{max_cycle}), \code{thresh} (\code{conv_tol}), \code{shift}
(\code{level_shift}) and \code{ene_tot} (\code{e_tot}).

\subsection{Results and methods}

After \code{kernel()}: \code{converged}, \code{e_tot}, \code{e_corr}
(relative to KRHF), \code{e_ref} (HF energy of the OBMP2 determinant), and
\code{mo_energy}, \code{mo_coeff} and \code{mo_occ} for all $k$-points of the
mesh (lists over $k$-points, occupied orbitals first). \code{fock_hf} holds
the HF Fock matrices in the basis of the final orbitals.

\begin{longtable}{@{}p{0.42\textwidth} p{0.52\textwidth}@{}}
\toprule
Method & Returns \\
\midrule
\endhead
\bottomrule
\endfoot
\code{get_fock()} & correlated Fock matrices (AO basis) at all $k$-points \\
\code{ip_koopmans()}, \code{ea_koopmans()} & minus the valence band maximum and minus the conduction band minimum on the mesh \\
\code{band_gap()} & gap of the eigenvalues of $\fbar$ on the mesh \\
\code{init_from_scf()} & takes the orbitals of the mean field without running OBMP2 (for conventional second-order quasiparticle energies) \\
\code{ipea(nvir=4, kpts_index=None)} & second-order IP/EA corrections; returns the quasiparticle gap \\
\code{self_energy(nvir=2, ...)} & tabulated second-order self-energy of the band edges \\
\code{quasiparticle_gap(scheme, reference)} & (direct gap at the first $k$-point, smallest gap on the mesh) \\
\code{peom_mp2(kpt_v=0, kpts_c=None)} & band edges and gaps from partitioned EOM-MP2 \\
\code{downfold(ncas, nelecas, cas_list=None)} & \code{KOBDF} object \\
\code{density_fit(auxbasis=None, with_df=None)} & sets another density-fitting object \\
\end{longtable}

\subsection{Restart}

To restart from given orbitals (for example a converged calculation with other
scaling factors), set \code{mp.mo_coeff}, \code{mp.mo_energy} and
\code{mp.mo_occ} (lists over the $k$-points of the mesh) before calling
\code{kernel()}.

\section{Algorithm}
\label{sec:kpt-algorithm}

\paragraph{Three-index integrals.} All integrals are products of GDF tensors,
\[
  (ia|jb) = \sum_Q L^Q_{ia}(k_i,k_a)\,L^Q_{jb}(k_j,k_b).
\]
The
occupied--virtual tensors (\emph{rows}, one per $k$-point) are transformed
once per cycle and kept in memory, or in a temporary HDF5 file when they need
more than half of \code{max_memory}. The AO tensors are cached between cycles
as far as the remaining memory allows (up to 80\,\% of \code{max_memory} in
total).

\paragraph{No integral with a general index.} The terms
$\sum \bar t_{iajb}(ia|pb)$ and $\sum \bar t_{iajb}(ia|jp)$ are evaluated
through the intermediate
\begin{equation}
  N^Q_{jb}(k_b) = \sum_{k_i}\sum_{ia} L^Q_{ia}(k_i,k_a)\,\bar t_{iajb},
\end{equation}
accumulated over $k_i$ and contracted once per pair $(k_j,k_b)$ with
$L^Q_{pb}$ and $L^Q_{jp}$, which are built on the fly.

\paragraph{Memory.} Nothing of size $\Nk^2 o^2 v^2$ is stored; the amplitudes
exist for one $k$-point quadruple at a time.

\paragraph{Symmetry.} Every term of the potential is written with its free
indices on $k_j$ (\emph{$j$ terms}) or on $k_b$ (\emph{$b$ terms}), using
$t_{iajb}=t_{jbia}$. $k_j$ then runs over the irreducible $k$-points only:
the $j$ terms are complete, and the $b$ terms are completed by the group
average $v = |G|^{-1}\sum_g g^{-1} P g$ of the partial sum $P$, with the
orbital rotations of \code{pyscf.pbc.symm}. The orbitals are kept at the
irreducible $k$-points and expanded in a symmetry-consistent gauge in every
cycle. If the mean-field solution does not have the symmetry of the mesh, the
full mesh is used. Only symmorphic operations are used (the default of
\pyscf{}).

\paragraph{Self-consistency.} The loop is \code{pycmf.core.scf_loop} with the
irreducible $k$-points as channels: DIIS with the commutator $[F,D]$,
convergence on $\Delta E$ and on the occupied--virtual block of $\fbar$, and
diagonalisation in the space of the current orbitals, so that orbitals
removed because of linear dependencies stay removed.

\section{Resources and performance}
\label{sec:kpt-resources}

\paragraph{Disk.} The three-index file of \pyscf{} holds all pairs of
$k$-points. For diamond gth-dzvp it takes 3.9~GB at $4\times4\times4$ and
14.8~GB at $5\times5\times5$, and \pyscf{} needs about as much again for a
temporary file while it is built. Set \code{mf.with_df._cderi_to_save} to
keep the file and \code{mf.with_df._cderi} to reuse it in a later run.

\paragraph{Memory.} The reading of the AO three-index integrals disappears
when \code{max_memory} is large enough to keep all AO tensors in memory (7.4~GB
for diamond gth-dzvp $4\times4\times4$; the default of 4000~MB holds a third of
them).

\paragraph{Timings.} Table~\ref{tab:ktimings} gives the time per
self-consistent cycle for diamond (two atoms per cell, GTH pseudopotential,
GDF, two threads of a shared virtual machine; timings fluctuate by about
$\pm20$\,\%). The script \code{benchmarks/kpoints_timing.py} reproduces such
measurements on other hardware.

\begin{table}[!htb]
\centering
\small
\caption{Time per KOBMP2 cycle for diamond, including the HF Fock build.}
\label{tab:ktimings}
\begin{tabular}{@{}l c r r r r@{}}
\toprule
Basis, mesh & $k$-points (irred.) & 0.2 & 1.0 full & 1.0 symm. & Memory \\
\midrule
gth-szv, $2^3$  & 8 (3)  & 4.2 s & 0.30 s & 0.17 s & 0.2 GB \\
gth-szv, $3^3$  & 27 (4) & 149 s & 3.7 s  & 2.1 s  & 0.3 GB \\
gth-dzvp, $2^3$ & 8 (3)  & 11 s  & 1.1 s  & 0.8 s  & 0.3 GB \\
gth-szv, $4^3$  & 64 (8) & $\approx$36 min & 33 s & 11.5 s & 0.9 GB \\
gth-dzvp, $3^3$ & 27 (4) & $\approx$7 min  & 27 s & 7.5 s  & 1.7 GB \\
gth-dzvp, $4^3$ & 64 (8) & $\approx$94 min & --   & 88 s   & 3.5 GB \\
\bottomrule
\end{tabular}
\end{table}

The values for version 0.2 have its debugging overhead removed; values marked
$\approx$ are extrapolated with $\Nk^3$. With DIIS, five to seven cycles bring
$|\Delta E|$ below $10^{-8}$~Hartree from HF orbitals; without DIIS seven to
thirteen.

\section{Nearly linearly dependent basis sets}
\label{sec:lindep}

The OBMP2 iterations are less robust than Hartree--Fock when the Bloch basis
is nearly linearly dependent. For diamond with gth-tzvp \pyscf{} removes two
to four functions per $k$-point, but overlap eigenvalues below $10^{-4}$
remain, and with them KOBMP2 does not converge. The function
\code{remove_linear_dependency_(mf, threshold=1e-4)}, called before
\code{mf.kernel()}, makes the mean field work in the space of the overlap
eigenvectors above the threshold:

\begin{lstlisting}[style=python]
from pycmf.pbc import KOBMP2, remove_linear_dependency_
mf = scf.KRHF(cell, kpts).density_fit()
remove_linear_dependency_(mf, 1e-4)
mf.kernel()
mp = KOBMP2(mf).run()
\end{lstlisting}

KOBMP2 keeps that orbital space and converges in five to six cycles, at a
cost of about $5\times10^{-5}$~Hartree in the HF energy. The number of
orbitals then depends on the $k$-point, the mean field is no longer exactly
symmetric, and KOBMP2 falls back to the full mesh.

\section{MPI}
\label{sec:kpt-mpi}

With \code{mpi4py} the loop over $k$-points of the correlation potential is
distributed. The same script runs on every process; rank 0 does the
Hartree--Fock part and broadcasts it, and the partial results are summed with
\code{allreduce}. Output is written by rank 0 only.

\begin{lstlisting}[style=shell]
$ OMP_NUM_THREADS=1 mpiexec -n 4 python kobmp2_mpi.py
\end{lstlisting}

See \code{examples/experimental/kobmp2_mpi.py}. Processes times threads must
not exceed the number of cores. For the small test case BN $3\times3\times3$
gth-szv, two processes reduce three cycles from 10.5~s to 8.2~s with symmetry;
the serial parts (HF build, reading of the integrals) limit the gain for
small cells.

\section{Band gaps and quasiparticle energies}
\label{sec:qp}

The eigenvalues of the correlated Fock matrix are not removal and addition
energies, so \code{band_gap()} is not a quasiparticle gap. Three levels of
theory for the band edges are built on a KOBMP2 calculation.

\subsection{Static correction: IPEA-KOBMP2}

\code{mp.ipea(nvir=4)} adds the second-order corrections
\begin{equation}
  \mathrm{IP}(h,k) = -\varepsilon(h,k) + 2\sum \bar t_{iajh}\,(ia|jh),\qquad
  \mathrm{EA}(l,k) = -\varepsilon(l,k) - 2\sum \bar t_{ialb}\,(ia|lb)
\end{equation}
to the converged eigenvalues (module \code{pbc/kipea.py}), stores
\code{mp.ip} and \code{mp.ea} for all $k$-points and returns the
quasiparticle gap $\min\mathrm{IP}-\max\mathrm{EA}$. The amplitudes carry
the same \code{c_os}, \code{c_ss} scaling as the potential. The corrections
are static and evaluated once; only the band edges are meaningful, because
the denominators can vanish for bands away from the gap.

\subsection{Second-order self-energy}

\code{pbc/kqp.py} tabulates the two parts of the diagonal second-order
self-energy on a frequency grid around each band energy,
\begin{equation}
  A_p(\omega) = \sum_{iab} \frac{W}{\omega + \varepsilon_i - \varepsilon_a - \varepsilon_b}
  \quad(\text{2p1h}),\qquad
  B_p(\omega) = \sum_{iaj} \frac{W}{\omega + \varepsilon_a - \varepsilon_i - \varepsilon_j}
  \quad(\text{2h1p}).
\end{equation}
With HF orbitals the diagonal of the OBMP2 potential is exactly one of them:
$v_{hh}=A_h(\varepsilon_h)$ for occupied and $v_{ll}=B_l(\varepsilon_l)$ for
virtual bands. The eigenvalues of $\fbar$ therefore contain the half of the
static self-energy that opens the gap ($\Sigma_{\mathrm{in}}$), and the IPEA
correction is the other half ($\Sigma_{\mathrm{out}}$). Nothing is counted
twice.

\begin{lstlisting}[style=python]
mp.kernel()
mp.self_energy(nvir=1, nocc=1)          # band edges only
direct, smallest = mp.quasiparticle_gap('dynamic')
\end{lstlisting}

\begin{table}[!htb]
\centering
\small
\caption{Quasiparticle schemes (\code{reference='obmp2'}).}
\begin{tabular}{@{}l l@{}}
\toprule
Scheme & Quasiparticle energy \\
\midrule
\code{'static'}  & $\varepsilon_p + \Sigma_{\mathrm{out}}(\varepsilon_p)$; identical to \code{mp.ipea()} \\
\code{'linear'}  & $\varepsilon_p + Z\,\Sigma_{\mathrm{out}}(\varepsilon_p)$, $Z = 1/(1 - A' - B')$ \\
\code{'dynamic'} & root of $\omega = \varepsilon_p + [\Sigma_{\mathrm{in}}(\omega) - \Sigma_{\mathrm{in}}(\varepsilon_p)] + \Sigma_{\mathrm{out}}(\omega)$ \\
\bottomrule
\end{tabular}
\end{table}

With \code{reference='hf'} together with \code{KOBMP2(mf).init_from_scf()}
the conventional second-order quasiparticle energies on Hartree--Fock are
obtained, $\varepsilon_p + A + B$.

Arguments of \code{self_energy}: \code{nvir} and \code{nocc} are the numbers
of lowest virtual and highest occupied bands that are treated
(\code{nocc=1, nvir=1} is enough for the gap); \code{kpts_index} selects
$k$-points (default: the irreducible ones); \code{grid} gives the frequency
offsets (default 41 points within $\pm0.4$~Hartree or 1.4 times the gap);
\code{scaling} is \code{'linear'} or \code{'hylleraas'}.

\paragraph{Spin-component scaling.} The amplitudes enter the potential
linearly with the factor $c$ and quadratically with $c^2$, so the half inside
the eigenvalues carries $2c-c^2$ per spin component, whereas the IPEA formula
multiplies the other half by $c$ itself. \code{self_energy(scaling='hylleraas')}
evaluates the self-energy with $2c-c^2$, so that both halves are weighted
alike; the default \code{'linear'} reproduces the IPEA numbers. For
$c_{\mathrm{OS}}=c_{\mathrm{SS}}=1$ there is no difference.

The lower-level functions are available directly:
\code{kqp.self_energy(mp, ...)}, \code{kqp.quasiparticle(se, scheme,
reference)}, \code{kqp.band_edges(se, e_qp)} and \code{kqp.gaps(se, e_qp)}.

\subsection{Partitioned EOM-MP2}

\code{pbc/kpeom.py} computes the band edges with partitioned EOM-MP2: the
eigenvalues of $e^{-T_2}\hat He^{T_2}$ in the space of 1h and 2h1p (1p and
2p1h) configurations, with the doubles--doubles block replaced by
orbital-energy differences. Compared with the second-order self-energy, one of
the two integrals of every term is dressed by the amplitudes, which adds the
leading ring (screening), ladder and exchange terms of third order. The method
has no adjustable parameter.

\begin{lstlisting}[style=python]
mp = KOBMP2(mf).run()                 # or KOBMP2(mf).init_from_scf()
res = mp.peom_mp2(kpt_v=0, kpts_c=[0, kc])
res['ip'], res['ea'][kc], res['gap'][kc]     # Hartree
\end{lstlisting}

Orbitals and orbital energies are those of the KOBMP2 object; the singles
block and the couplings use the HF Fock matrix in that basis. Only the target
$k$-point is an external index, so no array with three $k$-point indices is
stored. With \code{symmetry=True} (default) the singles space is the
degenerate level of the band edge; the neglected coupling to the other bands
of the $k$-point changes the electron affinity of diamond by 0.01--0.03~eV.
The IP and EA agree with the EOM-KCCSD code of \pyscf{} (with $t_1=0$,
MP2-like $t_2$ and the same doubles--doubles block) to $10^{-9}$~Hartree.

\subsection{Comparing gaps}

Band gaps with $k$-point sampling depend strongly on the mesh and on the
treatment of the exchange divergence (\code{exxdiv} of the mean field). A
value on a single mesh is not a result. Series with \code{exxdiv=None} and
\code{exxdiv='ewald'} must be extrapolated, for example with
$a + bN_k^{-1/3} + cN_k^{-2/3}$, before they can be compared with experiment.
The script \code{benchmarks/diamond_gaps.py} computes the gaps of diamond for
HF and the variants above on one mesh.

\section{k-point downfolding: KOBDF}
\label{sec:kobdf}

\code{KOBDF(mp, ncas, nelecas, cas_list=None)}, or \code{mp.downfold(ncas,
nelecas)}, builds an active-space Hamiltonian from a \code{KOBMP2} object,
per unit cell,
\begin{equation}
  \hat H_{\mathrm{act}} = E_{\mathrm{core}}
  + \frac1{\Nk}\sum_{k}\sum_{tu}\bar h_{tu}(k)\,\hat a^\dagger_{tk}\hat a_{uk}
  + \frac1{2\Nk}\sum_{k_1k_2k_3}\sum_{tuvw}
  h^{(2)}_{tuvw}(k_1,k_2,k_3)\,
  \hat a^\dagger_{tk_1}\hat a^\dagger_{vk_3}\hat a_{wk_4}\hat a_{uk_2},
\end{equation}
with $\bar h(k)=h(k)+v^{\mathrm{HF}}_{\mathrm{core}}(k)+v[T](k)-v[T_{\mathrm{act}}](k)$
and $h^{(2)}=(tk_1,uk_2|vk_3,wk_4)/\Nk$. The cost is one KOBMP2 cycle plus
one cycle restricted to the active bands. With
\code{KOBMP2(mf).init_from_scf().downfold(...)} the downfolding is done at
second order on HF orbitals.

\begin{lstlisting}[style=python]
mp = KOBMP2(mf).run()
ham = mp.downfold(ncas=8, nelecas=8)
h1, e_core = ham.get_h1eff()          # (nkpts, ncas, ncas), constant per cell
ham.get_cderi(k1, k2)                 # three-index integrals of the active bands
ham.get_h2eff(k1, k2, k3)             # (t k1, u k2 | v k3, w k4) / Nk
ham.scf(); ham.mo_energy; ham.band_gap()      # mean field of the model
ham.self_energy(nocc=1, nvir=1)
ham.quasiparticle_gap('dynamic')      # second-order gap of the model
h1s, h2s, e0, nelec = ham.supercell() # real Hamiltonian of the supercell
ham.as_scf()                          # PySCF RHF object (CCSD, CASCI, DMRG)
ham.kernel()                          # FCI of the supercell, energy per cell
\end{lstlisting}

\begin{description}[style=nextline,leftmargin=3.2cm,labelwidth=3cm]
  \item[\code{ncas}] number of active bands per $k$-point.
  \item[\code{nelecas}] number of active electrons per unit cell (even).
  \item[\code{cas_list}] 0-based band indices, the same at every $k$-point.
        By default the \code{ncas} bands around the Fermi level.
  \item[\code{conv_tol}, \code{max_cycle}] convergence of \code{scf()}
        (defaults $10^{-10}$ and 50).
  \item[\code{fcisolver}] solver of \code{kernel()}.
\end{description}

\paragraph{Mean field and quasiparticles of the model.} \code{scf()} solves
Hartree--Fock for $\hat H_{\mathrm{act}}$; its bands contain the correlation
with the orbitals outside the active space statically, through
$v^{\mathrm{ext}}$. \code{self_energy()} adds the correlation within the
active space with its frequency dependence, and
\code{quasiparticle(scheme)} solves for the bands as in \code{kqp}.

\paragraph{Many-body solvers.} \code{supercell()} returns real integrals of
the Born--von K\'arm\'an supercell of $\Nk$ cells ($\Nk n_{\mathrm{cas}}$
orbitals). The real orbitals are obtained from time reversal and are not
localized. \code{as_scf()} wraps the integrals in a \pyscf{} \code{RHF}
object; its energies refer to the supercell. \code{kernel()} runs FCI and is
only feasible for small models.

\paragraph{Sum rule.} \code{energy_tot()} over the reference determinant
equals $E_{\mathrm{KOBMP2}}-E_c[T_{\mathrm{act}}]$; the second term is stored
as \code{e_corr_act}.

\paragraph{Madelung term.} With \code{exxdiv='ewald'} the constant contains
$-\xi N_{\mathrm{act}}/2$ for the active electrons, and \code{get_fock}
shifts the occupied levels by $-\xi$, as \pyscf{} does for KRHF; the
integrals do not contain the term.

From \code{examples/experimental/kobdf.py} (diamond, gth-szv,
$2\times2\times1$, \code{exxdiv=None}, six active bands):

\begin{lstlisting}[style=output]
E(KRHF) = -9.26630050   E(KOBMP2) = -9.41369394   gap of the KOBMP2 bands = 0.37745
h1eff: (4, 6, 6)   E_core = -12.34762518
sum rule: <H_act> - [E(KOBMP2) - E_c(T_act)] = 8.2e-13
E(HF of H_act) = -9.34556626   gap = 0.32196
second-order gap of H_act, static  = 0.30080
second-order gap of H_act, linear  = 0.30205
second-order gap of H_act, dynamic = 0.30205
supercell: 24 orbitals, 24 electrons; E(HF) = -9.34556626, E(CCSD) = -9.39969464 per cell
\end{lstlisting}

\section{Unrestricted k-point OBMP2: KUOBMP2}

\code{KUOBMP2} is the four-index code of version 0.2, moved into the package
largely unchanged; its interface is that of version 0.2 (\code{thresh},
\code{niter}, \code{shift}). For a closed-shell cell it does not reproduce
KOBMP2 (see Section~\ref{sec:known-issues}), it has the cost problems of the
old code, and its ionization analysis uses hard-coded band indices. It should
be used only for comparison with old results.

\chapter{Code Structure and Development}
\label{ch:architecture}

This chapter is written for developers who want to understand, modify or
extend \pycmf{}.

\section{Repository layout}

\begin{lstlisting}
src/pycmf/
  __init__.py   version and top-level exports (stable methods only)
  core/         self-consistent driver, DIIS, occupation schemes,
                amplitudes, MPI helpers, experimental marker
  mol/          molecules: integrals, correlation potential, OBMP2,
                O2BMP2, OBDH, O2BDH, OO-DH, RO-OBMP2, NOCI
  downfold/     one-body downfolding, FCIDUMP export
  properties/   densities, dipole moments, <S^2>, hyperfine couplings,
                transition dipoles
  embed/        projection embedding, concentric localization (experimental)
  pbc/          k-point OBMP2, O2BMP2, quasiparticles, OBDF (experimental)
tests/          pytest suite
examples/       runnable scripts with reference output
benchmarks/     timing and gap scripts; inputs of version 0.2
docs/           theory, usage, validation, status, migration, this manual
\end{lstlisting}

\begin{longtable}{@{}p{0.30\textwidth} p{0.64\textwidth}@{}}
\caption{Modules of the package.}\\
\toprule
Module & Content \\
\midrule
\endfirsthead
\toprule
Module & Content \\
\midrule
\endhead
\bottomrule
\endfoot
\code{core/scf_loop.py} & self-consistent driver shared by all methods \\
\code{core/amplitudes.py} & first-order doubles amplitudes as spin blocks; \code{restrict_to_active} \\
\code{core/occupation.py} & occupation schemes \code{aufbau}, \code{mom}, \code{imom} \\
\code{core/mpi.py} & optional mpi4py helpers (\code{split}, \code{bcast}, \code{allreduce}) \\
\code{core/status.py} & \code{ExperimentalWarning} and \code{experimental(name)} \\
\code{mol/integrals.py} & integral backends \code{IncoreInts} (exact) and \code{DFInts} \\
\code{mol/bch.py} & correlation potential: first- and second-order BCH terms \\
\code{mol/obmp2.py} & \code{OBMP2Base}, \code{ROBMP2}, \code{UOBMP2}, factories \code{OBMP2}, \code{O2BMP2} \\
\code{mol/obdh.py} & \code{ROBDH}, \code{UOBDH}, factories \code{OBDH}, \code{O2BDH}, \code{FUNCTIONALS} \\
\code{mol/roobmp2.py} & \code{ROOBMP2} (experimental) \\
\code{mol/oodh.py} & \code{UOODH}, factory \code{OODH} (experimental): its own \code{build} returns the exact orbital gradient as the ov block of $\fbar$ \\
\code{mol/excited.py} & MOM-UHF excited determinants, channel selection \\
\code{mol/noci.py} & \code{NOCI}, \code{Determinant}, \code{make_basis} (experimental) \\
\code{downfold/obdf.py} & \code{OBDF} \\
\code{properties/rdm.py} & densities and dipole moments \\
\code{properties/spin.py} & $\langle S^2\rangle$ \\
\code{properties/hfcc.py} & isotropic hyperfine couplings \\
\code{properties/transition.py} & generalized Slater--Condon rules, transition dipoles (experimental) \\
\code{embed/partition.py} & SPADE partition, AO indices of atoms \\
\code{embed/cl.py} & concentric localization \\
\code{embed/projection.py} & \code{EmbeddedOBMP2}, \code{EmbeddedOBDH} \\
\code{pbc/kobmp2.py} & \code{KOBMP2}, \code{KO2BMP2}, \code{remove_linear_dependency_} \\
\code{pbc/kintegrals.py} & three-index GDF tensors in memory or HDF5 (\code{KDFInts}) \\
\code{pbc/kbch.py} & correlation potential with $k$-points \\
\code{pbc/ksymm.py} & irreducible $k$-points, expansion, symmetrization (\code{KSymmetry}) \\
\code{pbc/kipea.py} & static IP/EA corrections \\
\code{pbc/kqp.py} & second-order self-energy and quasiparticle energies \\
\code{pbc/kpeom.py} & partitioned EOM-MP2 \\
\code{pbc/kobdf.py} & \code{KOBDF} \\
\code{pbc/kuobmp2.py} & \code{KUOBMP2} (code of version 0.2) \\
\end{longtable}

\section{One self-consistent driver}

\code{core.scf_loop.kernel(mp, mo_energy, mo_coeff, nocc)} iterates over a
list of \emph{channels}. A channel is a spin channel for molecules (one for a
restricted, two for an unrestricted reference) and an irreducible $k$-point
for periodic systems. In each cycle the driver

\begin{enumerate}
  \item calls \code{mp.build(mo_energy, mo_coeff, nocc)}, which returns an
        \code{Intermediates} object with at least \code{e_tot},
        \code{fock_ao}, \code{fock_mo} and \code{dm_chan};
  \item prints the energy, $\Delta E$ and $\|\bar f_{ia}\|$ and checks
        convergence (both $|\Delta E|<$ \code{conv_tol} and
        $\|\bar f_{ia}\|<$ \code{conv_tol_grad});
  \item applies DIIS to the AO Fock matrices with the error vector $[F,D]$
        projected on a fixed orthonormal basis;
  \item diagonalises each Fock matrix in the span of the current orbitals,
        so that truncated or linearly dependent orbital spaces are preserved;
  \item occupies the new orbitals with \code{occupation.select}.
\end{enumerate}

It returns \code{(converged, e_tot, mo_energy, mo_coeff, intermediates)}.
The overlap matrix is taken from \code{mp.get_ovlp()} if the method defines
it (one per $k$-point) and from the mean field otherwise.

\section{The build step of molecular OBMP2}

\code{OBMP2Base.build} works in this order:

\begin{enumerate}
  \item HF Fock matrix from the current density (\code{get_hf()} returns the
        input mean field, or its HF counterpart for a Kohn--Sham object);
  \item integrals through \code{ao2mo()}: \code{IncoreInts} or
        \code{DFInts}, depending on \code{with_df};
  \item amplitude blocks: \code{core.amplitudes.make_blocks}, which applies
        \code{c_os}, \code{c_ss} and \code{level_shift};
  \item correlation potential and constant:
        \code{mol.bch.correlation_potential}, controlled by \code{bch_order};
  \item \code{get_reference(dm, h1e, vhf, fock_hf, e_hf)}, which returns the
        reference Fock matrix, the reference energy and the weight of $v$.
\end{enumerate}

The correlated Fock matrix is then the reference Fock matrix plus the weighted
potential. \code{ROBMP2} and \code{UOBMP2} differ only in \code{_occ} (2 or
1 electrons per occupied orbital of a channel) and in \code{_pack} and
\code{_unpack}, which convert between the arrays of \pyscf{} and the lists of
channels.

\section{Variants as hooks}

Variants of the method are hooks on this structure, not new modules:

\begin{itemize}
  \item \code{mol/obdh.py}: a mixin overrides \code{get_reference} so that
        it returns the Kohn--Sham Fock matrix and energy and the weight
        $a_c$.
  \item \code{mol/roobmp2.py}: changes the orbital update.
  \item \code{embed/projection.py}: adds the embedding potential in
        \code{get_reference} and computes the subsystem energy.
  \item \code{downfold/obdf.py}: calls \code{bch.correlation_potential}
        twice, with all amplitudes and with \code{restrict_to_active}.
  \item \code{pbc/kobmp2.py}: provides its own \code{build} and
        \code{get_ovlp} and reuses the driver with $k$-points as channels.
\end{itemize}

\section{Adding a method}

\begin{enumerate}
  \item Put the module at its final place (\code{mol/}, \code{pbc/} or
        \code{core/}) and call \code{pycmf.core.status.experimental(name)}
        in \code{__init__}. Do not export it in \code{pycmf/__init__.py}.
  \item Reuse \code{core.scf_loop.kernel} by providing
        \code{build(mo_energy, mo_coeff, nocc)}; derive from an existing
        class and override \code{get_reference} where possible.
  \item Add three kinds of tests: an exact limit (for example ``first cycle
        from HF orbitals equals \pyscf{} MP2'', ``all orbitals active equals
        FCI'' or ``$a_c=0$ is hybrid DFT''), a regression value, and input
        checks.
  \item When the method has been validated against published or independent
        results, remove the experimental marker, export it, and update
        \code{docs/status.md}.
\end{enumerate}

\section{Coding conventions}

\begin{itemize}
  \item \textbf{One method per module.} DF, DIIS, MOM, spin scaling and BCH
        order are attributes. Never add modules named \code{*_slow},
        \code{*_fast}, \code{*_mom} or \code{*_new}.
  \item \textbf{Placement.} Molecular code in \code{mol/}, periodic code in
        \code{pbc/}, shared code in \code{core/}. Reference implementations
        written as explicit loops belong in \code{tests/}.
  \item \textbf{\pyscf{} style.} Classes derive from
        \code{lib.StreamObject}; options are attributes listed in
        \code{_keys}; \code{kernel()} returns the main result; output goes
        through \code{pyscf.lib.logger}, never \code{print}.
  \item \textbf{Never modify the mean-field object} passed by the user. A test
        checks this.
  \item \textbf{Lint.} \code{ruff check src/pycmf tests examples} must pass
        (line length 110, single quotes). \code{src/pycmf/pbc} and
        \code{benchmarks} are excluded from ruff. The pre-commit hook runs
        \code{ruff format}.
\end{itemize}

\section{Development workflow}

\begin{lstlisting}[style=shell]
$ python3 -m venv .venv
$ env -u PYTHONPATH .venv/bin/pip install -e ".[dev]"
$ env -u PYTHONPATH .venv/bin/python -m pytest -m "not slow"   # before merging
$ env -u PYTHONPATH .venv/bin/python -m pytest                 # with slow tests
$ env -u PYTHONPATH .venv/bin/python -m pytest tests/test_obmp2.py::test_name
$ env -u PYTHONPATH .venv/bin/ruff check src/pycmf tests examples
$ cd examples && env -u PYTHONPATH ../.venv/bin/python check_examples.py --all
\end{lstlisting}

\code{pytest.ini} sets \code{pythonpath = src}, so the tests also run without
installing the package. The continuous integration
(\code{.github/workflows/ci.yml}) runs ruff and the fast tests on Python 3.9
and 3.12. Development happens on short-lived \code{feature/<name>} branches
that are rebased on the main branch; releases are tags.

\subsection*{Keeping the documentation consistent}

If a change alters behaviour, numbers or test counts, update the matching
document: \code{docs/theory.md} (equations), \code{docs/usage.md} (API),
\code{docs/status.md} (maturity and known issues),
\code{docs/validation.md} (what each test checks), \code{docs/kpoints.md}
($k$-point algorithm and timings), \code{docs/migration.md} (names of
version 0.2), and this manual (\code{docs/manual}).

\subsection*{Building this manual}

\begin{lstlisting}[style=shell]
$ cd docs/manual
$ latexmk -pdf pycmf_manual.tex       # or: pdflatex pycmf_manual (twice)
\end{lstlisting}

The manual uses only standard \LaTeX{} packages (\code{lmodern},
\code{booktabs}, \code{tabularx}, \code{longtable}, \code{listings},
\code{enumitem}, \code{fancyhdr}, \code{titlesec}, \code{hyperref}).

\chapter{Validation, Performance and Known Issues}
\label{ch:validation}

\section{Maturity of the methods}

\begin{longtable}{@{}p{0.24\textwidth} p{0.34\textwidth} p{0.34\textwidth}@{}}
\caption{What has been checked and what is open.}\\
\toprule
Method & Checked against & Open points \\
\midrule
\endfirsthead
\toprule
Method & Checked against & Open points \\
\midrule
\endhead
\bottomrule
\endfoot
OBMP2 (R, U; exact and DF) & MP2 at the first cycle; 0.2 energies to $10^{-10}$ & no frozen core, no analytic gradients \\
O2BMP2 & 0.2 energies & no $O(N^4)$ Laplace algorithm; no recommended \code{c_os} \\
MOM, IMOM & 0.2 energies & open-shell singlets are spin-mixed \\
OBDF & FCI and OBMP2 limits; 0.2 energies & FCIDUMP for restricted references only \\
OBDH, O2BDH & OBMP2 and hybrid-DFT limits; 0.2 energies & published parameters, not re-optimised; no benchmark \\
RO-OBMP2 (exp.) & 0.2 energies to $10^{-10}$ & no formal derivation, no DIIS, no benchmark \\
Embedding (exp.) & 0.2 energies; whole system and DFT-in-DFT limits & unrestricted only; no benchmark \\
OO-DH (exp.) & orbital gradient = finite differences ($10^{-8}$); ORCA 4.1 OO-RI-MP2 ($6\times10^{-8}$) and OO-RI-B2PLYP (H$_2$O, NH$_2$; lowering to $10^{-8}$); hybrid-DFT limit; relaxed dipole = finite field & unrestricted only; four-index blocks in memory; no nuclear gradients \\
Transition dipoles (exp.) & full-CI expansion of the two determinants ($10^{-9}$); unrelaxed single excitation; OB-O2B-simple scripts ($10^{-7}$) & determinant-level moment; $\sqrt2$ singlet approximation \\
NOCI (exp.) & $S$, $H$, $\mu$, $S^2$ against full CI ($10^{-9}$), also for a singular overlap; OB-O2B-simple scripts ($10^{-13}$ for the matrices) & effective Hamiltonians are model choices (issue 8) \\
KOBMP2, KO2BMP2 (exp.) & 0.2 code cycle by cycle; KMP2 at the first cycle; molecular code at $\Gamma$ & no published benchmark yet \\
KUOBMP2 (exp.) & -- & known inconsistency (issue 1) \\
KOBDF (exp.) & four-index KROBDF code; molecular OBDF at $\Gamma$; complete and empty active space; sum rule & closed shell; no Wannier90, band or DOS tools \\
\end{longtable}

``0.2 energies'' are total energies of the previous implementation (commit
29efd04) on small molecules. They are regression values, not an independent
validation. None of the published benchmark numbers has been recomputed with
version 1.0 yet.

\section{What the test suite checks}

\begin{longtable}{@{}p{0.55\textwidth} p{0.20\textwidth} p{0.17\textwidth}@{}}
\toprule
Property & Test file & Reference \\
\midrule
\endfirsthead
\toprule
Property & Test file & Reference \\
\midrule
\endhead
\bottomrule
\endfoot
First cycle from HF orbitals is MP2 (R, U) and DF-MP2 & \code{test_obmp2.py} & \code{pyscf.mp} \\
First-cycle correlated density is the unrelaxed MP2 density & \code{test_properties.py} & \code{pyscf.mp} \\
Restricted equals unrestricted for closed shells & all & internal \\
Converged orbital energies are eigenvalues of $\fbar$; $\bar f_{ia}=0$ & \code{test_obmp2.py} & internal \\
DIIS on and off give the same energy & \code{test_obmp2.py} & internal \\
The mean-field object is not modified & \code{test_obmp2.py} & internal \\
$c_{\mathrm{OS}}=c_{\mathrm{SS}}=0$ is HF; $c_{\mathrm{OS}}=c_{\mathrm{SS}}=1$ is OBMP2 & \code{test_o2bmp2.py} & internal \\
OBDH with $a_x=a_c=1$ is OBMP2; with $a_c=0$ hybrid DFT & \code{test_obdh.py} & \code{pyscf.dft} \\
OBDF with all orbitals active is FCI; with none it is OBMP2 & \code{test_obdf.py} & \code{pyscf.fci} \\
FCIDUMP round trip & \code{test_obdf.py} & \code{pyscf.tools} \\
Total energies of version 0.2 for all stable methods & all & version 0.2 \\
$\langle S^2\rangle$ and hyperfine couplings of OH & \code{test_properties.py} & version 0.2 \\
RO-OBMP2 and embedding; embedding limits & \code{test_experimental.py} & 0.2, \code{pyscf.dft} \\
Named double hybrids equal their four parameters & \code{test_obdh.py} & internal \\
OO-DH: orbital gradient = finite differences; OO-MP2 and OO-B2PLYP against ORCA; $a_c=0$ is hybrid DFT; relaxed dipole = finite field & \code{test_oodh.py} & ORCA 4.1, internal \\
Transition dipole: L\"owdin pair in the full CI space; $\langle 0|\mu|i\to a\rangle=-\langle i|r|a\rangle$; OB-O2B-simple scripts & \code{test_transition.py} & \code{pyscf.fci}, scripts \\
NOCI matrices against full CI (exact and direct integrals); diagonals; linear dependency; OB-O2B-simple scripts & \code{test_noci.py} & \code{pyscf.fci}, scripts \\
KOBDF: KROBDF code, $\Gamma$ point, limits, sum rule, supercell solvers & \code{test_kobdf.py} & KROBDF, \code{pyscf.fci}, \code{pyscf.cc} \\
KOBMP2: 0.2 code cycle by cycle, KMP2, $\Gamma$ point, symmetry, $k$-dependent occupations, disk storage, MPI & \code{test_kobmp2.py} & 0.2, \code{pyscf.pbc.mp} \\
\end{longtable}

Regression values use a tolerance of about $5\times10^{-9}$~Hartree. The
fixture \code{mf_cache} in \code{tests/conftest.py} provides shared
molecules for the whole session.

\section{Agreement with version 0.2}

\begin{table}[!htb]
\centering
\small
\caption{Largest deviation of total energies from version 0.2 (6-31G unless noted).}
\begin{tabularx}{\textwidth}{@{}X l r@{}}
\toprule
Method & Systems & $\max|\Delta E|$ (Ha) \\
\midrule
OBMP2, O2BMP2 (R, U, exact and DF, BCH 1 and 2, level shift) & H$_2$O, OH, CH$_2$ & $10^{-10}$ \\
state-specific OBMP2 / O2BMP2 (MOM) & H$_2$O & $2\times10^{-10}$ \\
OBDF, five active spaces & H$_2$/cc-pVDZ, LiH, H$_2$O & $4\times10^{-10}$ \\
OBDH, O2BDH & H$_2$O, CH$_2$, NH$_2$ & $4\times10^{-11}$ \\
OBDH, O2BDH & OH ($^2\Pi$, degenerate) & $10^{-7}$ \\
RO-OBMP2 & OH, CH$_2$, NH$_2$ & $10^{-13}$ \\
embedded OBMP2 / OBDH, with and without CL & (H$_2$O)$_2$ & $9\times10^{-11}$ \\
KOBMP2 & diamond, BN; up to 27 $k$-points & $3\times10^{-10}$ \\
\bottomrule
\end{tabularx}
\end{table}

\section{Performance of the molecular code}

Wall time per cycle for benzene/cc-pVDZ (114 basis functions, 21 occupied
orbitals) on two threads of a shared virtual machine (variations of about
20\,\% between runs; measure on your own hardware with
\code{benchmarks/timing.py}):

\begin{table}[!htb]
\centering
\small
\begin{tabular}{@{}l r r r@{}}
\toprule
Calculation & 0.2 (s) & 1.0 (s) & Ratio \\
\midrule
restricted, exact integrals     & 29.5 & 3.0 & 10 \\
unrestricted, exact integrals   & 13.4 & 7.6 & 1.8 \\
restricted, density fitting     & 10.5 & 1.8 & 6 \\
unrestricted, density fitting   & 7.6  & 3.6 & 2.1 \\
\bottomrule
\end{tabular}
\end{table}

The remaining cost with exact integrals is the $O(N^5)$ transformation of the
second index pair; use density fitting for larger molecules.

\section{Examples}

Every script in \code{examples/} has its expected output in
\code{examples/reference_output/}, and \code{check_examples.py} compares all
printed numbers with it.

\begin{table}[!htb]
\centering
\small
\begin{tabularx}{\textwidth}{@{}l X@{}}
\toprule
Script & Shows \\
\midrule
\code{01-obmp2.py} & restricted and unrestricted OBMP2, exact and density-fitted integrals \\
\code{02-o2bmp2.py} & O2BMP2 and the role of \code{c_os}: singlet--triplet gap of CH$_2$ \\
\code{03-excited_state.py} & state-specific excited state with MOM \\
\code{04-obdf.py} & downfolding along the N$_2$ dissociation curve, FCIDUMP export \\
\code{05-obdh.py} & OBDH and O2BDH, double hybrids by name or by parameters, open shells \\
\code{06-properties.py} & densities, dipole moment, $\langle S^2\rangle$, hyperfine couplings of CN \\
\code{experimental/ro_obmp2.py} & RO-OBMP2 for OH \\
\code{experimental/embedding.py} & OBDH and OBMP2 embedded in PBE, water dimer \\
\code{experimental/kobmp2.py} & KOBMP2 and KO2BMP2 for diamond, with and without symmetry \\
\code{experimental/kobmp2_mpi.py} & KOBMP2 with MPI \\
\code{experimental/kobdf.py} & $k$-point downfolding, quasiparticles and supercell solvers \\
\code{experimental/oodh.py} & OO-DH for several double hybrids, OO-MP2, relaxed dipole \\
\code{experimental/transition_dipole.py} & transition dipole and oscillator strength, OBMP2 and O2BMP2 \\
\code{experimental/noci.py} & NOCI of OBMP2 determinants with the three Hamiltonians \\
\bottomrule
\end{tabularx}
\end{table}

\section{Known issues}
\label{sec:known-issues}

\begin{enumerate}
  \item \textbf{KUOBMP2 does not reproduce KOBMP2 for a closed-shell cell.}
        For diamond (gth-szv, two $k$-points) KOBMP2 gives
        $-10.6408$~Hartree and KUOBMP2 $-8.6381$~Hartree from the same HF
        energy ($-10.5311$). The moved code reproduces the 0.2 numbers
        exactly, so the discrepancy is in the method code.
  \item \textbf{Hard-coded band indices in \code{kuobmp2.py}.} The
        ionization analysis addresses orbitals $n_{\mathrm{occ}}-3,\ldots,
        n_{\mathrm{occ}}+1$; cells with fewer than three occupied or two
        virtual bands per $k$-point raise \code{IndexError}.
  \item \textbf{KUOBMP2 is still the four-index code of 0.2} and has the cost
        problems that KOBMP2 had.
  \item \textbf{Embedding functional.} In 1.0 the two levels have separate
        functionals. The two conventions give very different energies
        ($-152.5729$ against $-152.6107$~Hartree for the water dimer of the
        tests); \code{xc=xc_env} reproduces version 0.2.
  \item \textbf{KROBDF with \code{exxdiv='ewald'}.} The constant of the old
        KROBDF code is higher by $\xi N_{\mathrm{act}}$ than the value for
        which the sum rule holds. \code{KOBDF} puts $-\xi N_{\mathrm{act}}/2$
        in the constant and the shift in \code{get_fock}.
  \item \textbf{Degenerate open shells with OBDH} converge only to
        $\|\bar f_{ia}\|\sim10^{-6}$ (Section~\ref{sec:obdh}).
  \item \textbf{The energy is not variational} (all methods except OO-DH),
        so its error is of first order in $\bar f_{ia}$. Use a tight
        \code{conv_tol_grad} when energy differences below $10^{-6}$~Hartree
        matter.
  \item \textbf{NOCI with \code{'state_specific_fock'} does not couple the
        $\alpha$ and $\beta$ branches of a channel.} They differ by two spin
        orbitals and the averaged Fock operator is a one-body operator. The
        pairs stay nearly degenerate, singlet and triplet are not separated,
        and $\mu$, $f$ and $\langle S^2\rangle$ of a pair depend on an
        arbitrary rotation (water/cc-pVDZ, 4$\to$5: two roots at 8.0759~eV,
        $\langle S^2\rangle\approx1$). \code{'bare'} and
        \code{'overlap_symmetric'} give the triplet at 6.72/7.73~eV and the
        singlet at 7.42/8.42~eV. The original OB-O2B-simple scripts show the
        same degeneracy.
\end{enumerate}

\section{Limitations}

\begin{itemize}
  \item No frozen core for molecules and no analytic gradients.
  \item No $O(N^4)$ Laplace-transformed algorithm for O2BMP2 or KO2BMP2.
  \item Periodic code: closed-shell \code{KRHF} references only (except the
        old KUOBMP2), three-dimensional systems, Gaussian density fitting,
        no GPU path, no real arithmetic for $\Gamma$-only meshes, no
        finite-size correction.
  \item P-EOM-MP2 keeps only the degenerate frontier level in the singles
        space and is not MPI-parallel.
  \item KOBDF: closed shell only; the supercell orbitals are not localized.
\end{itemize}


\appendix
\chapter{API Reference}
\label{ch:api}

This appendix lists the public classes and functions with their signatures.
Attributes and results are described in Chapters~\ref{ch:molecular} and
\ref{ch:periodic}. The docstrings in the source code are the authoritative
reference; \code{help(pycmf.OBMP2)} prints them.

\newcommand{\apientry}[2]{\par\noindent\code{#1}\par\vspace{-0.3em}\begin{list}{}{\leftmargin=1.5em\topsep=0pt}\item #2\end{list}\vspace{0.3em}}

\section{Top level: \texttt{pycmf}}

\code{pycmf.__version__} is \code{'1.0.0'}. Exported names:
\code{OBMP2}, \code{O2BMP2}, \code{ROBMP2}, \code{UOBMP2}, \code{OBDH},
\code{O2BDH}, \code{ROBDH}, \code{UOBDH}, \code{OBDF}, \code{properties}.

\apientry{OBMP2(mf, frozen=None, mo_coeff=None, mo_occ=None, **kwargs)}{%
Returns \code{UOBMP2} for a UHF or ROHF reference and \code{ROBMP2}
otherwise. \code{mo_coeff} and \code{mo_occ} replace the starting orbitals and
occupations of \code{mf}. \code{frozen} must be \code{None} or 0. Keyword
arguments are set as attributes.}

\apientry{O2BMP2(mf, c_os=None, frozen=None, mo_coeff=None, mo_occ=None, **kwargs)}{%
OBMP2 with \code{c_ss=0} and the given \code{c_os}. Raises
\code{ValueError} without \code{c_os}.}

\apientry{OBDH(mf, ax=None, ac=None, frozen=None, mo_coeff=None, mo_occ=None, xc=None, **kwargs)}{%
Returns \code{UOBDH} or \code{ROBDH}. \code{xc} is a name of
\code{FUNCTIONALS} or a \pyscf{} functional string; explicit \code{ax},
\code{ac} override the named parameters. \code{ax} together with a functional
string raises \code{ValueError}.}

\apientry{O2BDH(mf, c_os=None, ax=None, ac=None, frozen=None, mo_coeff=None, mo_occ=None, xc=None, **kwargs)}{%
OBDH with \code{c_ss=0} and the given \code{c_os} (mandatory).}

\apientry{OBDF(mp, ncas, nelecas, cas_list=None)}{%
Active-space Hamiltonian from a converged OBMP2 or O2BMP2 object. Methods:
\code{build()}, \code{get_h1eff()} $\to$ \code{(h1eff, e_core)},
\code{get_h2eff()}, \code{kernel(ci0=None, **kwargs)} $\to$ total energy,
\code{to_fcidump(filename, tol=1e-15, **kwargs)}.}

\section{Molecular classes: \texttt{pycmf.mol}}

\apientry{class OBMP2Base(lib.StreamObject)}{%
Base class. Methods: \code{kernel()}, \code{run()}, \code{dump_flags()},
\code{density_fit(auxbasis=None, with_df=None)}, \code{get_hf()},
\code{ao2mo()}, \code{build(mo_energy, mo_coeff, nocc)},
\code{get_reference(dm, h1e, vhf, fock_hf, e_hf)},
\code{make_amplitudes(mo_energy, mo_coeff, nocc, ints=None)},
\code{get_intermediates(with_amplitudes=False)},
\code{get_fock(ao_repr=False)}, \code{get_vcorr()},
\code{make_rdm1(ao_repr=True, correlated=False)},
\code{dip_moment(correlated=False, unit='Debye', verbose=None)},
\code{ip_koopmans()}, \code{ea_koopmans()}.}

\apientry{class ROBMP2(OBMP2Base)}{%
Closed-shell restricted reference. \code{spin_square()} returns
\code{(0, 1)}. Raises \code{TypeError} for UHF or ROHF objects.}

\apientry{class UOBMP2(OBMP2Base)}{%
Unrestricted reference (ROHF and RHF objects are converted). Additional
methods \code{spin_square(correlated=False)} and
\code{hfcc(atoms=None, correlated=False)}.}

\apientry{class Intermediates}{%
Quantities of one cycle: \code{e_tot}, \code{e_ref}, \code{e_corr},
\code{e_hf}, \code{dm_chan}, \code{fock_ao}, \code{fock_mo},
\code{fock_hf}, \code{v}, \code{c0}, \code{rdm}, \code{blocks},
\code{nocc}.}

\apientry{class ROBDH(ROBMP2), class UOBDH(UOBMP2)}{%
Double hybrids. Additional methods \code{set_functional(name)},
\code{get_xc()}, \code{get_ks()}.}

\apientry{obdh.FUNCTIONALS, obdh.double_hybrid(name)}{%
Dictionary name $\to$ \code{(ax, ac, xc_x, xc_c)} with the keys
\code{B2PLYP}, \code{MPW2PLYP}, \code{B2TPLYP}, \code{B2GPPLYP},
\code{B2KPLYP}, \code{PBE0DH}, \code{PBEQIDH}, \code{PBE02};
\code{double_hybrid} normalises a name (case, hyphens, underscores) and returns
the tuple or \code{None}.}

\apientry{class roobmp2.ROOBMP2(UOBMP2)}{%
\code{ROOBMP2(mf, frozen=None, mo_coeff=None, mo_occ=None)}; ROHF reference,
experimental.}

\section{Properties: \texttt{pycmf.properties}}

\apientry{rdm.make_rdm1(mp, ao_repr=True, correlated=False)}{One-particle density matrix.}
\apientry{rdm.dip_moment(mp, correlated=False, unit='Debye', verbose=None)}{Dipole moment.}
\apientry{spin.spin_square(mp, correlated=False)}{$\langle S^2\rangle$ and $2S+1$ (unrestricted).}
\apientry{hfcc.isotropic(mp, atoms=None, correlated=False)}{Isotropic hyperfine coupling constants in MHz.}

\section{Embedding: \texttt{pycmf.embed} (experimental)}

\apientry{EmbeddedOBMP2(mf, active_atoms=None, **kwargs)}{OBMP2 embedded in DFT.}
\apientry{EmbeddedOBDH(mf, active_atoms=None, **kwargs)}{OBDH embedded in DFT.}
\apientry{spade_partition(mol, s1e, mo_occ, atoms, return_sigma=False)}{SPADE partition of the occupied orbitals.}
\apientry{atom_aos(mol, atoms)}{AO indices of the given atoms.}
\apientry{concentric_localization(mo_vir, s1e, fock, active_aos, n_shells=1, tol=1e-5, n_span=None)}{Virtual orbitals of the first CL shells.}
\apientry{shell_sizes(mo_vir, s1e, fock, active_aos, tol=1e-5)}{Sizes of all CL shells.}

\section{Periodic systems: \texttt{pycmf.pbc} (experimental)}

\apientry{KOBMP2(mf, frozen=None, kpt_symmetry=None)}{%
Methods: \code{kernel()}, \code{run()}, \code{density_fit(auxbasis=None,
with_df=None)}, \code{get_ovlp()}, \code{build(mo_energy, mo_coeff, nocc)},
\code{get_fock()}, \code{ip_koopmans()}, \code{ea_koopmans()},
\code{band_gap()}, \code{init_from_scf()},
\code{ipea(nvir=4, kpts_index=None)},
\code{self_energy(nvir=2, kpts_index=None, grid=None, mo_energy=None, nocc=None, scaling='linear')},
\code{quasiparticle_gap(scheme='static', reference='obmp2')},
\code{peom_mp2(kpt_v=0, kpts_c=None, symmetry=True)},
\code{downfold(ncas, nelecas, cas_list=None)}.}

\apientry{KO2BMP2(mf, c_os=None, **kwargs)}{%
\code{KOBMP2} with \code{c_ss=0}; \code{c_os} is mandatory; further keyword
arguments (\code{frozen}, \code{kpt_symmetry}) go to the constructor.}

\apientry{remove_linear_dependency_(mf, threshold=1e-4)}{%
Canonical orthogonalisation per $k$-point for a mean-field object, before
\code{mf.kernel()}. Modifies and returns \code{mf}.}

\apientry{KOBDF(mp, ncas, nelecas, cas_list=None)}{%
Methods: \code{build()}, \code{get_h1eff()}, \code{get_cderi(k1, k2)},
\code{get_h2eff(k1=None, k2=None, k3=None)}, \code{get_jk(dm)},
\code{reference_rdm1()}, \code{make_rdm1(mo_rot=None)},
\code{get_fock(dm=None)}, \code{energy_tot(dm=None)},
\code{scf(dm0=None)}, \code{band_gap()},
\code{self_energy(nocc=None, nvir=None, kpts_index=None, grid=None, c_os=1.0, c_ss=1.0)},
\code{quasiparticle(scheme='dynamic')},
\code{quasiparticle_gap(scheme='dynamic')}, \code{supercell(tol=1e-6)},
\code{as_scf()}, \code{kernel(ci0=None, **kwargs)}.}

\apientry{kqp.self_energy(mp, nvir=2, kpts_index=None, grid=None, mo_energy=None, nocc=None, scaling='linear')}{%
List over the $k$-points of dictionaries with \code{'occ'}, \code{'e'},
\code{'grid'}, \code{'A'}, \code{'B'}, \code{'A0'}, \code{'B0'},
\code{'dA'}, \code{'dB'} (\code{None} for $k$-points not treated).}
\apientry{kqp.quasiparticle(se, scheme='static', reference='obmp2', e_ref=None)}{%
\code{(e_qp, z)}, lists over the $k$-points.}
\apientry{kqp.band_edges(se, e_qp)}{Quasiparticle energies of the valence and conduction band edges per $k$-point.}
\apientry{kqp.gaps(se, e_qp)}{%
(direct gap at the first $k$-point, smallest gap, $k$-index of the valence
band maximum, $k$-index of the conduction band minimum).}
\apientry{kpeom.gap(mp, kpt_v=0, kpts_c=None, symmetry=True)}{%
Dictionary with \code{'ip'}, \code{'ea'} and \code{'gap'}.}
\apientry{kpeom.ip(mp, kpt=0, nbands=None, symmetry=True, grid=None), kpeom.ea(...)}{%
P-EOM-MP2 ionization potential and electron affinity at one $k$-point.}
\apientry{kipea.kernel(mp, nvir=4, kpts_index=None)}{%
\code{(ip, ea)}, lists over the $k$-points.}

\apientry{KUOBMP2(mf, frozen=0, mo_coeff=None, mo_occ=None)}{%
Unrestricted $k$-point OBMP2 of version 0.2 with its old interface.}

\section{Core: \texttt{pycmf.core}}

\apientry{scf_loop.kernel(mp, mo_energy, mo_coeff, nocc, verbose=None)}{%
Self-consistent driver; returns \code{(converged, e_tot, mo_energy,
mo_coeff, intermediates)}.}
\apientry{amplitudes.make_blocks(ints, mo_energy, nocc, c_os=1.0, c_ss=1.0, shift=0.0)}{%
List of spin blocks \code{(x, y, t, tbar)}.}
\apientry{amplitudes.restrict_to_active(blocks, occ_act, vir_act)}{Blocks with all indices in the active space.}
\apientry{occupation.select(mo_energy, mo_coeff, nocc, scheme='aufbau', ref_occ=None, s1e=None)}{%
Reorders the orbitals so that the selected ones come first.
\code{occupation.SCHEMES = ('aufbau', 'mom', 'imom')}.}
\apientry{mpi.rank, mpi.size, mpi.split(tasks), mpi.bcast(obj), mpi.allreduce(arrays)}{%
MPI helpers; no-ops without \code{mpi4py} or with \code{PYCMF_NO_MPI} set.}
\apientry{status.ExperimentalWarning, status.experimental(name, note='')}{%
Warning class and the function that emits it.}

\chapter{Migrating from pyCMF 0.2}
\label{ch:migration}

Version 1.0 is a rewrite on a common core. One class per method replaces the
many modules of version 0.2, and variants are attributes.

\begin{longtable}{@{}p{0.45\textwidth} p{0.49\textwidth}@{}}
\caption{Names of version 0.2 and their replacement.}\\
\toprule
Version 0.2 & Version 1.0 \\
\midrule
\endfirsthead
\toprule
Version 0.2 & Version 1.0 \\
\midrule
\endhead
\bottomrule
\endfoot
\code{OBMP2(mf)}, \code{OBMP2_slow}, \code{OBMP2_einsum} & \code{pycmf.OBMP2(mf)} with an RHF reference \\
\code{UOBMP2(mf)}, \code{UOBMP2_slow} & \code{pycmf.OBMP2(mf)} with a UHF or ROHF reference \\
\code{DFOBMP2}, \code{DFUOBMP2}, \code{DFUOBMP2_einsum} & \code{pycmf.OBMP2(mf_df)} or \code{.density_fit()} \\
\code{UOBMP2_SCS}; \code{css}, \code{cos} & \code{c_ss}, \code{c_os}; \code{pycmf.O2BMP2(mf, c_os=...)} \\
\code{UOBMP2_mom}, \code{UOBMP2_mom_diis}, \code{mom_select=True} & \code{occupation='imom'} \\
\code{UHF_mom_diis} & \code{pyscf.scf.addons.mom_occ} \\
\code{second_order=True/False}, \code{ampf} & \code{bch_order=2} / \code{1} \\
\code{thresh}, \code{niter}, \code{shift} & \code{conv_tol}, \code{max_cycle}, \code{level_shift} \\
\code{mp.ene_tot} & \code{mp.e_tot} \\
\code{kernel()} returning a tuple & \code{kernel()} returns \code{e_tot}; the rest are attributes \\
\code{OBMP2_downfold}, \code{UOBMP2_downfold} with intermediates copied by hand & \code{pycmf.OBDF(mp, ncas, nelecas, cas_list=...)} \\
\code{h1mo_act_eff}, \code{h2mo_act}, \code{ene_inact} & \code{get_h1eff()} $\to$ \code{(h1eff, e_core)}, \code{get_h2eff()} \\
\code{caslist} counted from 1 & \code{cas_list} counted from 0 \\
\code{DFTOBMP2}, \code{OBDH_CL} with \code{alphaa=(ax, ac)} & \code{pycmf.OBDH(mf, ax=..., ac=...)} or \code{xc='B2PLYP'} \\
\code{OBDH_CL}, \code{OBMP2_CL} with \code{use_embed=True} & \code{pycmf.embed.EmbeddedOBDH}, \code{EmbeddedOBMP2} \\
\code{make_fc(dm)} & \code{mp.hfcc(correlated=False)} \\
OB-O2B-simple: \code{calc_ob_o2b_transition.py} & \code{transition.transition_dipole(gs, es, singlet=True)} \\
OB-O2B-simple: \code{calc_noci_ob_o2b.py}, \code{run_mom_hf}, \code{select_channels} & \code{noci.make_basis}, \code{noci.NOCI}, \code{excited.mom_excitation}, \code{excited.select_channels} \\
\code{oo_dh.py}: \code{UOODH(mf, xc).kernel()}, \code{dh_energy} & \code{pycmf.OODH(mol, xc=...).run()}; \code{e_dh} is the non-optimized double hybrid \\
\code{make_rdm1(use_t2=...)} & \code{mp.make_rdm1(correlated=...)} \\
\code{ROOBMP2} (\code{OBMP/ro_obmp2.py}) & \code{pycmf.mol.roobmp2.ROOBMP2} \\
\code{KOBMP2} ($k$-point, restricted) & \code{pycmf.pbc.KOBMP2}; old names still work \\
\code{KUOBMP2} & \code{pycmf.pbc.KUOBMP2} (interface unchanged) \\
\code{KROBDF} with intermediates copied by hand & \code{mp.downfold(ncas, nelecas)} $\to$ \code{KOBDF} \\
\code{make_fock_df(scf_active=True)}, \code{make_ipea()} & \code{ham.scf()}, \code{ham.self_energy()} and \code{ham.quasiparticle()} \\
\code{Supercell}, \code{solve_fci}, \code{solve_ccsd} & \code{ham.supercell()}, \code{ham.kernel()}, \code{ham.as_scf()} \\
\end{longtable}

\section*{Example}

Version 0.2:
\begin{lstlisting}[style=python]
mp = UOBMP2_mom_diis(mf)
mp.second_order = True; mp.css = 0.0; mp.cos = 1.2
mp.mom_select = True; mp.thresh = 1e-8
e = mp.kernel()
\end{lstlisting}

Version 1.0:
\begin{lstlisting}[style=python]
mp = pycmf.O2BMP2(mf, c_os=1.2, occupation='imom')
e = mp.kernel()
\end{lstlisting}

\section*{Changes in behaviour}

\begin{itemize}
  \item \code{bch_order=2} is the default everywhere. In 0.2 \code{OBMP2}
        defaulted to the first-order variant and \code{DFOBMP2} to second
        order.
  \item Spin scaling acts on the amplitudes, as in the O2BMP2 paper. The
        different scaling of \code{uobmp2_scs.py} was not carried over.
  \item The mean-field object is no longer modified by a calculation.
  \item Convergence is tested on the energy and on $\bar f_{ia}$.
  \item DIIS uses the commutator $[F,D]$ per spin, with at most
        \code{diis_space} vectors.
  \item \code{frozen} orbitals raise \code{NotImplementedError} for molecules.
  \item Not carried over: the second-order Koopmans correction
        \code{make_IPEA} for molecules, the CAS variants, the
        \code{stabilize_scf} helper (use \code{mf.stability()}), and the
        Wannier90, real-space Hamiltonian, band-structure and DOS tools of the
        KROBDF package.
  \item The input scripts of version 0.2 are kept in
        \code{benchmarks/legacy-0.2/}. They use the old interface and must be
        ported before use.
\end{itemize}

\chapter{References}
\label{ch:references}

\subsection*{Methods implemented in pyCMF}

\begin{enumerate}[label={[\arabic*]}]
  \item T. N. Lan and T. Yanai, Correlated one-body potential from
        second-order M\o ller--Plesset perturbation theory: Alternative to
        orbital-optimized MP2 method, \emph{J. Chem. Phys.} \textbf{138},
        224108 (2013).
  \item L. N. Tran, Improving perturbation theory for open-shell molecules
        via self-consistency, \emph{J. Phys. Chem. A} \textbf{125}, 9242
        (2021).
  \item L. N. Tran, Can second-order perturbation theory accurately predict
        electron density of open-shell molecules? The importance of
        self-consistency, \emph{Phys. Chem. Chem. Phys.} \textbf{24}, 19393
        (2022).
  \item N. T. Le and L. N. Tran, \emph{J. Phys. Chem. A} \textbf{127}, 5222
        (2023). [one-body downfolding]
  \item N. T. Tran, H. T. Nguyen and L. N. Tran, \emph{J. Phys. Chem. A}
        \textbf{128}, 1543 (2024). [O2BMP2]
  \item N. T. Tran and L. N. Tran, \emph{J. Chem. Phys.} \textbf{162}, 104109
        (2025). [state-specific excited states]
\end{enumerate}

\subsection*{Software}

\begin{enumerate}[label={[\arabic*]},resume]
  \item Q. Sun \emph{et al.}, PySCF: the Python-based simulations of
        chemistry framework, \emph{WIREs Comput. Mol. Sci.} \textbf{8}, e1340
        (2018).
\end{enumerate}

\subsection*{Other methods used}

\begin{enumerate}[label={[\arabic*]},resume]
  \item S. Grimme, \emph{J. Chem. Phys.} \textbf{124}, 034108 (2006). [B2PLYP]
  \item T. Schwabe and S. Grimme, \emph{Phys. Chem. Chem. Phys.} \textbf{8},
        4398 (2006). [mPW2PLYP]
  \item A. Tarnopolsky \emph{et al.}, \emph{J. Phys. Chem. A} \textbf{112}, 3
        (2008). [B2T-PLYP, B2K-PLYP]
  \item A. Karton \emph{et al.}, \emph{J. Phys. Chem. A} \textbf{112}, 12868
        (2008). [B2GP-PLYP]
  \item E. Br\'emond and C. Adamo, \emph{J. Chem. Phys.} \textbf{135}, 024106
        (2011). [PBE0-DH]
  \item E. Br\'emond \emph{et al.}, \emph{J. Chem. Phys.} \textbf{141},
        031101 (2014). [PBE-QIDH]
  \item J.-D. Chai and S.-P. Mao, \emph{Chem. Phys. Lett.} \textbf{538}, 121
        (2012). [PBE0-2]
  \item A. T. B. Gilbert, N. A. Besley and P. M. W. Gill, \emph{J. Phys.
        Chem. A} \textbf{112}, 13164 (2008). [MOM]
  \item G. M. J. Barca, A. T. B. Gilbert and P. M. W. Gill, \emph{J. Chem.
        Theory Comput.} \textbf{14}, 1501 (2018). [IMOM]
  \item F. R. Manby \emph{et al.}, \emph{J. Chem. Theory Comput.} \textbf{8},
        2564 (2012). [projection embedding]
  \item D. Claudino and N. J. Mayhall, \emph{J. Chem. Theory Comput.}
        \textbf{15}, 1053 (2019). [SPADE]
  \item D. Claudino and N. J. Mayhall, \emph{J. Chem. Theory Comput.}
        \textbf{15}, 6085 (2019). [concentric localization]
  \item M. Nooijen and J. G. Snijders, \emph{J. Chem. Phys.} \textbf{102},
        1681 (1995). [partitioned EOM]
  \item M. F. Lange and T. C. Berkelbach, \emph{J. Chem. Phys.} \textbf{155},
        081101 (2021). [P-EOM-MP2 for solids]
  \item U. Bozkaya, J. M. Turney, Y. Yamaguchi, H. F. Schaefer and
        C. D. Sherrill, \emph{J. Chem. Phys.} \textbf{135}, 104103 (2011).
        [OO-MP2]
  \item F. Neese, T. Schwabe, S. Kossmann, B. Schirmer and S. Grimme,
        \emph{J. Chem. Theory Comput.} \textbf{5}, 3060 (2009).
        [orbital-optimized MP2 and double hybrids, OO-RI-B2PLYP]
  \item R. Peverati and M. Head-Gordon, \emph{J. Chem. Phys.} \textbf{139},
        024110 (2013). [orbital-optimized double hybrids]
  \item S. Bourne Worster, O. Feighan and F. R. Manby, \emph{J. Chem. Phys.}
        \textbf{154}, 124106 (2021). [transition properties of
        non-orthogonal determinants]
  \item Y. Lemke, J. Kussmann and C. Ochsenfeld, \emph{J. Chem. Theory
        Comput.} \textbf{21}, 10193 (2025). [generalized Slater--Condon rules,
        NOCI]
\end{enumerate}


\end{document}
